% !TeX program = pdflatex
% Compile with pdfLaTeX until cross-references and the contents are stable.

\documentclass[11pt,a4paper]{article}
\usepackage[T1]{fontenc}
\usepackage[utf8]{inputenc}
\usepackage{lmodern,microtype}
\usepackage{geometry}
\geometry{left=23mm,right=23mm,top=22mm,bottom=24mm,headheight=22pt}
\usepackage{setspace}
\setstretch{1.06}
\usepackage{amsmath,amssymb,mathtools,bm}
\usepackage{array,tabularx,booktabs}
\usepackage{enumitem,xcolor,tcolorbox}
\tcbuselibrary{skins,theorems,breakable}
\usepackage{needspace,etoolbox,fancyhdr,titlesec}
\usepackage{hyperref}

\definecolor{ink}{HTML}{414141}
\definecolor{vuBurgundy}{HTML}{78003F}
\definecolor{blue}{HTML}{174F78}
\definecolor{navy}{HTML}{123653}
\definecolor{paleBlue}{HTML}{F1F6F9}
\definecolor{paleBurgundy}{HTML}{FAF2F6}
\definecolor{workGray}{HTML}{ABB3BB}
\definecolor{softGray}{HTML}{F5F5F5}
\color{ink}
\setlength{\parindent}{0pt}
\setlength{\parskip}{0.48em}
\setlength{\emergencystretch}{1.6em}
\raggedbottom
\clubpenalty=10000
\widowpenalty=10000
\displaywidowpenalty=10000
\setlist[enumerate]{leftmargin=1.8em,itemsep=0.25em,topsep=0.25em}
\setlist[itemize]{leftmargin=1.5em,itemsep=0.25em,topsep=0.25em}
\setcounter{tocdepth}{2}
\setcounter{secnumdepth}{2}
\numberwithin{equation}{section}
\titleformat{\section}{\Large\bfseries\color{vuBurgundy}}{\thesection}{0.65em}{}
\titleformat{\subsection}{\large\bfseries\color{blue}}{\thesubsection}{0.65em}{}
\titlespacing*{\section}{0pt}{1.5em}{0.55em}
\titlespacing*{\subsection}{0pt}{1.1em}{0.4em}
\pretocmd{\section}{\Needspace{12\baselineskip}}{}{}
\pretocmd{\subsection}{\Needspace{5\baselineskip}}{}{}
\pagestyle{fancy}
\fancyhf{}
\fancyhead[L]{\small\color{vuBurgundy}\textbf{General Relativity I}}
\fancyhead[R]{\small\color{ink!65}Lecture notes}
\fancyfoot[C]{\small\color{vuBurgundy}\thepage}
\renewcommand{\headrulewidth}{0.4pt}
\renewcommand{\headrule}{\hbox to\headwidth{\color{vuBurgundy}\leaders\hrule height\headrulewidth\hfill}\vskip-\headrulewidth}
\hypersetup{pdftitle={General Relativity I - Lecture Notes},
  pdfsubject={Manifolds, tangent spaces, differential forms, connections and metrics},
  colorlinks=true,linkcolor=navy,urlcolor=blue,bookmarksnumbered=true}

\newcommand{\R}{\mathbb R}
\newcommand{\dd}{\mathrm d}
\newcommand{\Cinf}{C^\infty}
\newcommand{\X}{\mathfrak X}
\newcommand{\Lie}{\mathcal L}
\newcommand{\id}{\operatorname{id}}
\newcommand{\grad}{\operatorname{grad}}
\newcommand{\GL}{\mathrm{GL}}
\newcommand{\slideref}[1]{\par\nobreak\smallskip\nobreak{\footnotesize\color{ink!60}\textsc{Source slide #1}}\par\nobreak\smallskip\nobreak}
\newcommand{\slidesref}[1]{\par\nobreak\smallskip\nobreak{\footnotesize\color{ink!60}\textsc{Source slides #1}}\par\nobreak\smallskip\nobreak}
\newcommand{\proofend}{\unskip\nobreak\hfill\(\square\)\par}
\newenvironment{proof}{\par\noindent\textit{Proof.}\ }{\proofend}

\tcbset{statement/.style={enhanced,breakable,lines before break=4,boxrule=0.65pt,arc=1.2mm,
  left=2.5mm,right=2.5mm,top=1.6mm,bottom=1.6mm,
  before skip=7pt,after skip=8pt,fonttitle=\bfseries,coltitle=white}}
\newtcbtheorem[number within=section]{definition}{Definition}{statement,unbreakable,
  title after break={Definition~\thetcbcounter\ (continued)},
  colback=paleBlue,colframe=blue,colbacktitle=blue}{def}
\newtcbtheorem[use counter from=definition]{theorem}{Theorem}{statement,
  title after break={Theorem~\thetcbcounter\ (continued)},
  colback=paleBurgundy,colframe=vuBurgundy,colbacktitle=vuBurgundy}{thm}
\newtcbtheorem[use counter from=definition]{proposition}{Proposition}{statement,
  title after break={Proposition~\thetcbcounter\ (continued)},
  colback=paleBurgundy,colframe=vuBurgundy,colbacktitle=vuBurgundy}{prop}
\newtcbtheorem[use counter from=definition]{examplebox}{Example}{statement,
  title after break={Example~\thetcbcounter\ (continued)},
  colback=softGray,colframe=ink!45,colbacktitle=ink!65}{ex}
\newtcolorbox{remarkbox}[1][Remark]{statement,title={#1},title after break={#1 (continued)},
  colback=paleBurgundy,colframe=vuBurgundy,colbacktitle=vuBurgundy}
% Blank writing space that can continue across a page break.
\newcommand{\workingspace}[1]{%
  \par
  \begingroup
  \offinterlineskip
  \dimen0=#1\relax
  \loop\ifdim\dimen0>0pt
    \dimen2=5mm
    \ifdim\dimen0<\dimen2
      \dimen2=\dimen0
    \fi
    \hbox to \linewidth{%
      \vrule width 0pt height \dimen2 depth 0pt
      \hfil
    }%
    \penalty0
    \advance\dimen0 by -\dimen2
  \repeat
  \endgroup
}

% The optional argument specifies blank space AFTER the exercise prompt.
% Keep each exercise prompt and its writing space together on one page.
\newtcolorbox{studentwork}[2][3.3cm]{%
  enhanced,
  unbreakable,
  colback=white,
  colframe=workGray,
  colbacktitle=softGray,
  coltitle=ink,
  title={Student work -- #2},
  title after break={Student work -- #2 (continued)},
  fonttitle=\bfseries,
  boxrule=0.6pt,
  arc=1mm,
  left=3mm,
  right=3mm,
  top=2mm,
  bottom=2mm,
  before skip=8pt,
  after skip=11pt,
  after upper={\workingspace{#1}}
}

\begin{document}
\hypersetup{pageanchor=false}
\begin{titlepage}
\centering
\vspace*{2.2cm}
{\Huge\bfseries\color{vuBurgundy}General Relativity I\par}
\vspace{0.45cm}
{\LARGE\color{vuBurgundy}Lecture Notes\par}
\vspace{0.8cm}
{\large Manifolds, tangent spaces, differential forms,\\[0.25em]
connections, and Lorentzian geometry\par}
\vfill
{\small Metric signature \((+,-,-,-)\)\par}
\vspace{0.35cm}
{\small Smooth means \(C^\infty\)\par}
\end{titlepage}
\hypersetup{pageanchor=true}
\pagenumbering{roman}
\tableofcontents
\clearpage
\pagenumbering{arabic}

\section{From affine spacetime to local geometry}
\label{sec:opening}
\slideref{1}

In special relativity, spacetime is an affine space equipped with the
Minkowski metric. Given two events \(p\) and \(q\), their difference \(q-p\)
is a well-defined displacement vector. All such displacements belong to
the same vector space \(V\). A vector \(v\in V\) can describe a displacement
starting at any event, taking \(p\) to \(p+v\) or \(q\) to \(q+v\). The
arrows from \(p\) to \(p+v\) and from \(q\) to \(q+v\) therefore represent
the same displacement vector, even though they start at different events.

In general relativity, spacetime is described by a smooth four-dimensional
manifold \(M\) equipped with a Lorentzian metric \(g\). No global affine
structure is assumed. Each event \(p\) has a tangent space \(T_pM\).
A smooth parametrized curve passing through \(p\) determines a tangent
vector there, given by its derivative. The space \(T_pM\) consists of all
vectors obtained in this way. Vectors in \(T_pM\) can be added and
subtracted, but a vector at another event \(q\) belongs to a different
space, \(T_qM\). The manifold structure alone does not tell us which vector
in \(T_qM\) should correspond to a given vector in \(T_pM\). A connection
provides a rule for making this comparison by transporting vectors along
a chosen curve between the two events.

This generalization must be distinguished from a change of coordinates.
Changing coordinates assigns different numerical labels to the same events,
without changing the geometry. For example, a nonlinear coordinate change
can make the components of the Minkowski metric depend on position.
Returning to inertial Cartesian coordinates makes them constant again.
The position dependence in the other coordinates does not indicate
curvature.

General relativity allows both flat and curved spacetimes. If curvature
is nonzero at an event \(p\), no coordinates can make all metric components
constant throughout a neighborhood of \(p\). One can still arrange
\(g_{\mu\nu}(p)=\eta_{\mu\nu}\) at the single event, where
\(\eta_{\mu\nu}\) is the constant Minkowski matrix. This does not make
\(g_{\mu\nu}(x)=\eta_{\mu\nu}\) in a neighborhood. The relevant distinction
is therefore not whether the displayed components vary, but whether the
metric can be made constant throughout a coordinate neighborhood.

We begin by making local coordinates and tangent vectors precise. This
will let us describe fields and their derivatives before adding a metric.
The familiar linear algebra of SR1 remains valid within each tangent
space, while the new geometric structures explain how these spaces fit
together.

\subsection{Linear structure at a point}
\slidesref{2--4}

The vector-space and tensor conventions of SR1 apply at each point of a
manifold. Unless stated otherwise, scalars are real. Let \(V\) be an
\(n\)-dimensional real vector space. Its dual space is
\[
 V^*=\operatorname{Hom}_{\R}(V,\R),
\]
where \(\operatorname{Hom}_{\R}(V,W)\) denotes the real vector space of
linear maps from \(V\) to another real vector space \(W\). Elements of
\(V^*\) are covectors.

Let \(e_1,\ldots,e_n\) be a basis of \(V\). Its dual basis
\(\varepsilon^1,\ldots,\varepsilon^n\) satisfies
\(\varepsilon^i(e_j)=\delta^i{}_j\). Thus
\[
 v=v^i e_i,\qquad \omega=\omega_i\varepsilon^i,\qquad
 \omega(v)=\omega_i v^i.
\]
An index repeated exactly twice in a term, once upstairs and once
downstairs, is summed over its range. Latin coordinate indices run
from \(1\) to \(n\), and Greek spacetime indices run from \(0\) to
\(3\). The pairing \(\omega(v)\) requires no metric.

For nonnegative integers \(r,s\), a tensor of type \((r,s)\) belongs to
\[
 T^r_s(V)=V^{\otimes r}\otimes(V^*)^{\otimes s}.
\]
A zeroth tensor power is \(\R\), so a tensor of type \((0,0)\) is a
scalar. A tensor is a finite sum of tensor products. In the chosen
bases it has a unique component expansion
\[
 T=T^{i_1\cdots i_r}{}_{j_1\cdots j_s}\,
 e_{i_1}\otimes\cdots\otimes e_{i_r}\otimes
 \varepsilon^{j_1}\otimes\cdots\otimes\varepsilon^{j_s}.
\]
It acts multilinearly on \(r\) covectors and \(s\) vectors through
their dual pairings. In particular, a bilinear form has type \((0,2)\).
These constructions do not require a metric.

A change of basis changes components, not the tensor. If
\(e'_a=A^i{}_a e_i\), then
\(\varepsilon'^a=(A^{-1})^a{}_i\varepsilon^i\). For a tensor of type
\((1,1)\), this gives
\[
 \mathbf T'=A^{-1}\mathbf T A,
 \qquad \mathbf T=(T^i{}_j),\quad \mathbf T'=(T'^a{}_b).
\]
This is the change-of-basis rule from SR1 for an arbitrary invertible
matrix \(A\).

\begin{studentwork}[6cm]{basis-independent quantities}
Using the change-of-basis law for a tensor of type \((1,1)\),
prove that the trace and determinant of its component matrix
are independent of the chosen basis.
\end{studentwork}

On a manifold, take \(V=T_pM\) at each point \(p\). Local bases and
the matrices relating them may vary smoothly with \(p\).

\section{Smooth manifolds and local coordinates}
\label{sec:manifolds}
\slideref{5}

\subsection{Charts and smooth structures}

A coordinate system assigns \(n\) real numbers to each point in an open
region of an \(n\)-dimensional manifold. These numbers label the point
and need not be components of a position vector. A coordinate system
need not cover the whole manifold, so its domain must be specified.
Where two coordinate domains overlap, a transition map relates the
two coordinate descriptions of the same points.

\begin{definition}{Topological manifold and coordinate chart}{chart}
An \(n\)-dimensional \emph{topological manifold} is a Hausdorff,
second-countable topological space \(M\) in which every point lies in
an open set \(U\) admitting a homeomorphism
\[
 \varphi:U\longrightarrow\varphi(U)\subset\R^n
\]
onto an open subset of \(\R^n\). The pair \((U,\varphi)\) is a
\emph{coordinate chart}. Writing
\(\varphi(p)=(x^1(p),\ldots,x^n(p))\) defines the coordinate functions
\(x^i:U\to\R\).
\end{definition}

A homeomorphism is a continuous bijection with continuous inverse. The
Hausdorff condition requires distinct points to have disjoint open
neighborhoods. Second countability requires a countable collection of
open sets such that every open set is a union of members of that collection.
The existence of a Euclidean coordinate neighborhood around
every point does not by itself guarantee that the whole space is
Hausdorff or second countable. These two conditions are therefore
imposed separately in the definition of a manifold.

Let \((U,\varphi)\) and \((V,\psi)\) be coordinate charts with
\(U\cap V\neq\varnothing\). Every point \(p\in U\cap V\) has two
coordinate descriptions,
\[
 x=\varphi(p),
 \qquad
 y=\psi(p).
\]
To obtain \(y\) from \(x\), first apply \(\varphi^{-1}\) to recover
the point \(p\), then apply \(\psi\) to find its coordinates in the
second chart. The composition therefore acts as
\[
 x=\varphi(p)
 \xrightarrow{\;\varphi^{-1}\;}
 p
 \xrightarrow{\;\psi\;}
 \psi(p)=y.
\]
It changes the coordinates used to describe the point, not the point
itself. The resulting \emph{transition map} is
\[
 \psi\circ\varphi^{-1}:
 \varphi(U\cap V)\longrightarrow\psi(U\cap V).
\]
The restriction to the overlap is necessary because the point must
belong to the domains of both charts. Its inverse is
\(\varphi\circ\psi^{-1}\), which converts the second coordinate
description back to the first.

Since \(U\cap V\) is open and the charts are homeomorphisms onto open
subsets of \(\R^n\), both \(\varphi(U\cap V)\) and \(\psi(U\cap V)\)
are open in \(\R^n\). The transition maps therefore act between open
subsets of Euclidean space, so their smoothness can be tested using
ordinary multivariable calculus.

\begin{definition}{Smooth atlas and smooth manifold}{smooth-manifold}
Two charts are \emph{smoothly compatible} if the transition map
\(\psi\circ\varphi^{-1}\) is \(C^\infty\) on \(\varphi(U\cap V)\)
and its inverse is \(C^\infty\) on \(\psi(U\cap V)\).
Charts with disjoint domains are compatible by convention.
A \emph{smooth atlas} is a family of pairwise smoothly compatible
charts whose domains cover \(M\).

A smooth atlas is \emph{maximal} if it contains every chart that is
smoothly compatible with all of its charts. A \emph{smooth structure}
is a maximal smooth atlas. A topological manifold equipped with
a smooth structure is a \emph{smooth manifold}.
\end{definition}

Any smooth atlas determines a unique maximal one by adjoining every
chart compatible with all of its charts. To check compatibility of
two added charts near a point, compose their transitions through a
chart in the original atlas containing that point. Thus specifying
a convenient covering atlas suffices to specify the smooth structure.

\subsection{Smooth maps}

Charts let us test smoothness using ordinary calculus. For a map between
manifolds, choose coordinates on both the domain and the target, then
test the resulting map between open subsets of Euclidean spaces.

\begin{definition}{Smooth map and diffeomorphism}{smooth-map}
Let \(M\) and \(N\) be smooth manifolds. A map \(F:M\to N\) is
\emph{smooth} if, for every \(p\in M\), there are charts
\((U,\varphi)\) and \((V,\psi)\) from their respective smooth
structures, with \(p\in U\) and \(F(U)\subset V\), such that
\[
 \psi\circ F\circ\varphi^{-1}:
 \varphi(U)\longrightarrow\psi(V)
\]
is smooth in the ordinary Euclidean sense.

A \emph{diffeomorphism} is a smooth bijection whose inverse is smooth.
\end{definition}

This definition does not depend on the chosen charts. Replacing either
chart composes the coordinate representative with smooth transition
maps on the relevant overlaps. In particular, a scalar field
\(f:M\to\R\) is smooth precisely when \(f\circ\varphi^{-1}\) is smooth
in every chart belonging to the chosen smooth structure. Their space
is denoted by \(\Cinf(M)\). Pointwise addition, scalar multiplication,
and multiplication of functions make it a real algebra. This algebra
will provide the functions on which tangent vectors act as differentiation
rules.

Let \((U,\varphi)\) and \((V,\psi)\) be smooth charts on \(M\).
The coordinate representatives of a scalar field \(f:M\to\R\) are
\[
 f_x=f\circ\varphi^{-1},
 \qquad
 f_y=f\circ\psi^{-1},
\]
defined on \(\varphi(U)\) and \(\psi(V)\), respectively. For
\(p\in U\cap V\), write \(x=\varphi(p)\) and \(y=\psi(p)\). Then
\[
 f_y(y)=f_x(x)=f(p).
\]
Equivalently, on \(\psi(U\cap V)\),
\[
 f_y=f_x\circ(\varphi\circ\psi^{-1}).
\]
The coordinate expressions may differ, but they assign the same
value to the same point. This generalizes the scalar-field
transformation law of SR1 from affine Lorentz coordinate changes
to arbitrary smooth coordinate transitions.

\subsection{Local coordinates and global structure}
\slideref{6}

The definition of a manifold guarantees a coordinate chart around every
point. It does not guarantee a single chart whose domain is the whole
manifold. Such a \emph{global chart} would be a homeomorphism from \(M\)
onto an open subset of \(\R^n\).

Consider the unit sphere
\[
 S^2=\{(x,y,z)\in\R^3:x^2+y^2+z^2=1\}.
\]
It is compact because it is closed and bounded in \(\R^3\). If a global
chart \(\varphi:S^2\to W\) existed, its image \(W\) would be a nonempty
open subset of \(\R^2\). It would also be compact, since a homeomorphism
preserves compactness. No nonempty open subset of \(\R^2\) is compact,
so such a chart cannot exist. We can instead cover the sphere with two
charts, each omitting a different point.

\begin{examplebox}{Stereographic coordinates on the sphere}{stereographic}
Let \(N=(0,0,1)\) and \(S=(0,0,-1)\) be the north and south poles.
Define the open subsets
\[
 U_N=S^2\setminus\{N\},
 \qquad
 U_S=S^2\setminus\{S\}.
\]
The subscript indicates the pole excluded from the chart domain.

For \(p=(x,y,z)\in U_N\), stereographic projection from \(N\) assigns
to \(p\) the intersection of the line through \(N\) and \(p\) with
the plane \(z=0\). Identifying this plane with \(\R^2\) gives a pair
of coordinates. The line has the parametrization
\[
 N+t(p-N)=(tx,ty,1+t(z-1)).
\]
Its third coordinate vanishes when \(t=1/(1-z)\). Applying the same
construction from \(S\) gives the two maps
\(\varphi_N:U_N\to\R^2\) and \(\varphi_S:U_S\to\R^2\), with
\[
 \varphi_N(x,y,z)=\left(\frac{x}{1-z},\frac{y}{1-z}\right),
 \qquad
 \varphi_S(x,y,z)=\left(\frac{x}{1+z},\frac{y}{1+z}\right).
\]

Each map assigns two coordinates to every point in its domain.
Conversely, every pair \((u,v)\in\R^2\) determines exactly one point
in the corresponding domain. Writing \(q=u^2+v^2\), the inverse maps
are
\[
 \varphi_N^{-1}(u,v)=\frac{(2u,2v,q-1)}{1+q},
 \qquad
 \varphi_S^{-1}(u,v)=\frac{(2u,2v,1-q)}{1+q}.
\]
Both maps and their inverses are continuous, so they are
homeomorphisms onto \(\R^2\). In particular,
\(\varphi_N(S)=\varphi_S(N)=(0,0)\). Thus the origin in the
north-pole chart represents the south pole, while the origin in
the south-pole chart represents the north pole.

The two chart domains cover the whole sphere, and their overlap
excludes both poles,
\[
 U_N\cup U_S=S^2,
 \qquad
 U_N\cap U_S=S^2\setminus\{N,S\}.
\]
In either chart, the coordinate image of this overlap is
\(\R^2\setminus\{(0,0)\}\). The transition map is therefore
\[
 \varphi_S\circ\varphi_N^{-1}:
 \R^2\setminus\{(0,0)\}\longrightarrow\R^2\setminus\{(0,0)\},
\]
with
\[
 (\varphi_S\circ\varphi_N^{-1})(u,v)
 =\left(\frac{u}{u^2+v^2},\frac{v}{u^2+v^2}\right).
\]
The reverse transition \(\varphi_N\circ\varphi_S^{-1}\) has the same
formula. The denominator is nonzero throughout the stated domain,
so both transitions are \(C^\infty\). Consequently,
\((U_N,\varphi_N)\) and \((U_S,\varphi_S)\) form a smooth atlas on
\(S^2\).
\end{examplebox}

The exclusion of \((0,0)\) applies only to the transition map, not
to either full chart image. The south pole is covered by \(U_N\),
and the north pole is covered by \(U_S\). Neither pole belongs to
the overlap, because neither has coordinates in both charts.


\begin{studentwork}[3.6cm]{the transition between two sphere charts}
Verify the inverse and transition formulas. Compute the Jacobian
determinant of the transition map and show that it is nonzero
everywhere on its domain.
\end{studentwork}

\slideref{7}
The smooth structure specifies which functions and curves on the
manifold are smooth. Smooth transition maps ensure that this property
does not depend on the chart used to test it. The smooth structure
alone does not determine lengths, angles, causal cones, or parallel
transport. A metric supplies scalar products in tangent spaces and,
in the Lorentzian case, light cones. A connection supplies a rule for
parallel transport along curves.

We used the representation of \(S^2\) as a surface in \(\R^3\) to
construct its charts explicitly. The definition of a smooth manifold
does not require such a surrounding space. The following constructions
use smooth curves and functions on the manifold itself, without
requiring an embedding in a larger Euclidean space.


Unless stated otherwise, all manifolds below are finite-dimensional,
smooth, and without boundary. Smooth means \(C^\infty\). In particular,
all charts used for differentiation belong to the chosen smooth structure.

\section{Curves, tangent vectors, and derivations}
\label{sec:tangent}
\slideref{8}

\subsection{Differentiating along a curve}

A smooth parametrized curve is a smooth map \(c:I\to M\), where
\(I\subset\R\) is an open interval. The parameter \(t\in I\) specifies
which point \(c(t)\) on the curve is being considered. It need not represent
physical time. Suppose \(c(t_0)=p\). We want to define the tangent vector
\(\dot c(t_0)\), which describes how the curve passes through \(p\).

In an affine space, this vector can be obtained from a difference quotient
of points. On a general manifold, the subtraction
\(c(t_0+h)-c(t_0)\) is not defined. Instead, let \(f\) be a smooth real-valued
function near \(p\). The composition \(f\circ c\) is a real-valued function
of one variable, so its derivative is defined by ordinary calculus. We set
\begin{equation}
 \dot c(t_0)(f)
 :=\left.\frac{\dd}{\dd t}\right|_{t=t_0} f(c(t)).
 \label{eq:curve-action}
\end{equation}
Thus \(\dot c(t_0)\) is described by its action on smooth functions. Given
\(f\), it returns the rate at which \(f\) changes along the curve at
\(t_0\). It is the rule \(f\mapsto\dot c(t_0)(f)\), not any one of its
numerical values, that represents the tangent vector.

Ordinary differentiation is linear and obeys the product rule. Writing
\(v=\dot c(t_0)\), these properties become
\[
 v(af+bh)=a\,v(f)+b\,v(h),\qquad
 v(fh)=v(f)h(p)+f(p)v(h),
\]
where \(a,b\in\R\). The values \(f(p)\) and \(h(p)\) appear because the
derivative is evaluated at the parameter value for which \(c(t_0)=p\).
These properties will also define tangent vectors without first choosing a
curve.

\subsection{Local functions and tangent vectors}

Only the behavior of \(f\) near \(p\) matters in
\eqref{eq:curve-action}. To avoid requiring a function to be defined on all
of \(M\), consider pairs \((U,f)\), where \(U\) is an open neighborhood of
\(p\) and \(f\in\Cinf(U)\). Two pairs \((U,f)\) and \((W,h)\) are
identified when \(f=h\) on some open neighborhood of \(p\) contained in
\(U\cap W\). The resulting equivalence class is the \emph{germ of \(f\)
at \(p\)}, denoted by \([f]_p\). Equality of germs means agreement on a
neighborhood, not just equality of the two values at \(p\).

The set of smooth germs is denoted by \(\Cinf_p(M)\). Addition and
multiplication are performed after restricting representatives to a common
neighborhood. These operations do not depend on the representatives, so
\(\Cinf_p(M)\) is a real algebra. Evaluation \([f]_p\mapsto f(p)\) is
also well defined. We usually write \(v(f)\) instead of \(v([f]_p)\).

\begin{definition}{Tangent vector}{tangent-vector}
A \emph{tangent vector at \(p\)} is an \(\R\)-linear map
\[
 v:\Cinf_p(M)\longrightarrow\R
\]
satisfying the Leibniz rule
\[
 v(fh)=v(f)h(p)+f(p)v(h).
\]
Such a map is called a \emph{derivation at \(p\)}. The space of all
these derivations is the \emph{tangent space} \(T_pM\). Its vector-space
operations are defined by
\[
 (v+w)(f)=v(f)+w(f),\qquad (av)(f)=a\,v(f).
\]
\end{definition}

Equation~\eqref{eq:curve-action} therefore defines an element of \(T_pM\).
The coordinate description below will show that every element of \(T_pM\)
can be obtained from a curve. Every derivation annihilates constants.
Indeed, \(v(1)=v(1\cdot1)=2v(1)\), so \(v(1)=0\), and linearity then
gives \(v(a)=0\) for each constant \(a\in\R\).

\subsection{The coordinate basis}
\slidesref{9--10}

Let \((U,\varphi)\) be a smooth chart about \(p\), with coordinate
functions \(x^1,\ldots,x^n\). Write \(u^1,\ldots,u^n\) for the ordinary
Cartesian coordinates on the image \(\varphi(U)\subset\R^n\). Thus
\(x^i=u^i\circ\varphi\). A function \(f\) on \(U\) is represented in
these coordinates by \(\widetilde f=f\circ\varphi^{-1}\).

Differentiate this representative with respect to its \(i\)-th coordinate
and evaluate at \(\varphi(p)\). This defines an operator on functions,
\[
 \partial_i|_p(f)
 :=\left.
 \frac{\partial(f\circ\varphi^{-1})}{\partial u^i}
 \right|_{\varphi(p)},
 \qquad
 \partial_i=\frac{\partial}{\partial x^i}.
\]
The notation \(\partial_i f(p)\) abbreviates the same value. Since a
partial derivative is linear and satisfies the product rule,
\(\partial_i|_p\) is a tangent vector. Applying it to a coordinate
function gives
\[
 \partial_i|_p(x^j)=\delta_i{}^j.
\]

\begin{theorem}{Coordinate basis of the tangent space}{coordinate-basis}
The vectors \(\partial_1|_p,\ldots,\partial_n|_p\) form a basis of
\(T_pM\). Every tangent vector has a unique expansion
\begin{equation}
 v=v^i\partial_i|_p,\qquad v^i=v(x^i),\qquad
 v(f)=v^i\partial_i f(p).
 \label{eq:tangent-components}
\end{equation}
In particular, \(\dim T_pM=\dim M=n\).
\end{theorem}

\begin{proof}
If \(b^i\partial_i|_p=0\), applying both sides to \(x^j\) gives
\(b^j=0\). Hence the coordinate vectors are linearly independent.

To prove that they span \(T_pM\), we must show that a derivation is
completely determined by its values on the coordinate functions.
Let \(v\) be any derivation at \(p\) and put \(a=\varphi(p)\).
Restrict to a convex coordinate neighborhood of \(a\). The fundamental
theorem of calculus along the segment from \(a\) to \(u\) expresses
the change in \(f\) in terms of the coordinate increments,
\[
 \widetilde f(u)-\widetilde f(a)=(u^i-a^i)h_i(u),\qquad
 h_i(u)=\int_0^1
 \frac{\partial\widetilde f}{\partial u^i}(a+t(u-a))\,\dd t.
\]
The functions \(h_i\) are smooth, and \(h_i(a)=\partial_i f(p)\).
On the corresponding neighborhood of \(p\), this identity reads
\[
 f-f(p)=(x^i-a^i)(h_i\circ\varphi).
\]
Apply \(v\) and use the Leibniz rule. Constants are annihilated, and
\(x^i(p)-a^i=0\), so
\[
 \begin{aligned}
 v(f)
 &=v(x^i)h_i(a)+(x^i(p)-a^i)v(h_i\circ\varphi)\\
 &=v(x^i)\partial_i f(p).
 \end{aligned}
\]
This is the action of \(v(x^i)\partial_i|_p\) on every germ. Hence it is
\(v\). Uniqueness follows from linear independence.
\end{proof}

The integral identity is exact for smooth functions. It does not assume
that \(f\) is analytic or equal to a Taylor series. The theorem shows
that a derivation is determined by the \(n\) numbers \(v(x^i)\), although
it acts on an infinite-dimensional algebra of functions. Only first
partial derivatives of \(f\) enter its value.

\subsection{Recovering the velocity of a curve}

Write \(x^i(t)=x^i(c(t))\). The ordinary chain rule gives
\[
 \left.\frac{\dd}{\dd t}\right|_{t=t_0}f(c(t))
 =\dot x^i(t_0)\partial_i f(p),\qquad p=c(t_0).
\]
Comparing this with \eqref{eq:tangent-components} yields
\begin{equation}
 \dot c(t_0)=\dot x^i(t_0)\partial_i|_{c(t_0)}.
 \label{eq:velocity-components}
\end{equation}
Thus the derivatives \(\dot x^i(t_0)\) are the components of the tangent
vector in the coordinate basis. They are not the basis vectors themselves.

Conversely, let \(v=v^i\partial_i|_p\). For sufficiently small \(t\),
the curve
\[
 c(t)=\varphi^{-1}\bigl(\varphi(p)+t(v^1,\ldots,v^n)\bigr)
\]
is defined, passes through \(p\) at \(t=0\), and has \(\dot c(0)=v\).
The straight line used in this formula lies in the coordinate image.
Its inverse image need not look straight in an embedding of the manifold.
This construction proves that every tangent vector is a curve velocity.

Taking the component tuple to be the \(i\)-th standard vector of
\(\R^n\) gives a curve whose velocity is \(\partial_i|_p\). Along that
curve, only \(x^i\) changes, at unit parameter rate. Two curves through
\(p\) determine the same tangent vector precisely when their coordinate
velocities agree in one chart, and hence in every chart.

For the helix \(c(t)=(r\cos t,r\sin t,bt)\) in \(\R^3\), with fixed
\(r>0\) and \(b\in\R\), the usual velocity components are
\((-r\sin t,r\cos t,b)\). As a tangent derivation, the same velocity is
\[
 \dot c(t)=-r\sin t\,\partial_x|_{c(t)}
            +r\cos t\,\partial_y|_{c(t)}
            +b\,\partial_z|_{c(t)}.
\]
Its action on \(f\) is the corresponding linear combination of partial
derivatives of \(f\), evaluated at \(c(t)\).

Under a regular change of parameter \(\widetilde c(s)=c(t(s))\),
\[
 \widetilde c'(s_0)=t'(s_0)\,\dot c(t(s_0)).
\]
Regularity means that the parameter change is a local diffeomorphism, so
\(t'(s_0)\neq0\). For a nonzero velocity, a positive factor preserves
its oriented tangent direction and a negative factor reverses it.
Changing the parameter can therefore change the tangent vector even when
the curve traces the same set of points.

The velocity records the first-order behavior of a curve at the point.
The following notation lets us compare that behavior for two smooth curves.
For a vector-valued function \(h(t)\), the notation \(h(t)=O(t^2)\)
as \(t\to0\) means that \(\|h(t)\|\leq C|t|^2\) for some constant
\(C\) and all sufficiently small \(t\).

\begin{studentwork}[3.5cm]{first-order equivalence of curves}
Let \(c_1(0)=c_2(0)=p\) for two smooth curves. Prove that
\(\dot c_1(0)=\dot c_2(0)\) if and only if
\[
 \varphi(c_1(t))-\varphi(c_2(t))=O(t^2)
 \qquad (t\to0)
\]
in one chart about \(p\). Show that the condition then holds in every
chart. Give two parametrizations of the same curve with different
nonzero velocities at \(p\).
\end{studentwork}

\section{Differentials and coordinate transformations}
\label{sec:differentials}
\slidesref{9--10}

\subsection{The differential of a map}

Let \(F:M\to N\) be smooth. A curve \(c\) through \(p\) is sent to the
curve \(F\circ c\) through \(F(p)\). The \emph{differential} of \(F\)
relates their tangent vectors. Its input lies in \(T_pM\), and its output
lies in \(T_{F(p)}N\).

To define the output as a derivation, let \(f\) be a smooth function
near \(F(p)\). The composition \(f\circ F\) is a smooth function near
\(p\), so a vector \(v\in T_pM\) can act on it.

\begin{definition}{Differential or pushforward}{differential}
For a smooth map \(F:M\to N\), its \emph{differential at \(p\)} is the
linear map
\[
 \dd F_p:T_pM\longrightarrow T_{F(p)}N,
 \qquad (\dd F_p v)(f)=v(f\circ F).
\]
It is also called the \emph{pushforward at \(p\)}. The function \(f\)
represents a germ at \(F(p)\).
\end{definition}

If two representatives of \(f\) agree near \(F(p)\), their compositions
with \(F\) agree near \(p\). Thus the definition is independent of the
representative. The Leibniz rule gives
\[
 (\dd F_pv)(fh)
 =(\dd F_pv)(f)\,h(F(p))+f(F(p))\,(\dd F_pv)(h),
\]
so the output is a tangent vector at \(F(p)\). Linearity in \(v\)
follows directly from the definition.

When \(v=\dot c(0)\), its action is
\[
 (\dd F_pv)(f)
 =\left.\frac{\dd}{\dd t}\right|_{t=0}f(F(c(t))).
\]
This is the action of the velocity of \(F\circ c\). Hence
\[
 \dd F_p(\dot c(0))=(F\circ c)'(0).
\]

Choose coordinates \(x^i\) near \(p\) and \(y^a\) near \(F(p)\).
The functions \(F^a=y^a\circ F\) give the target coordinates of the
image point as functions on the domain. Applying the differential to
the coordinate basis gives
\begin{equation}
 \dd F_p(\partial_i|_p)
 =\frac{\partial F^a}{\partial x^i}(p)
   \left.\frac{\partial}{\partial y^a}\right|_{F(p)}.
 \label{eq:jacobian-differential}
\end{equation}
Indeed, the \(a\)-th component of the left-hand side is its value on
\(y^a\), namely \(\partial_i(y^a\circ F)(p)\). The Jacobian is therefore
the matrix of \(\dd F_p\) in these two coordinate bases. It need not be
square when the manifolds have different dimensions.

For smooth maps \(F:M\to N\) and \(G:N\to P\), the chain rule is
\[
 \dd(G\circ F)_p=\dd G_{F(p)}\circ\dd F_p.
\]
Both sides send \(v\) to the derivation
\(h\mapsto v(h\circ G\circ F)\). If \(F\) is a diffeomorphism, applying
this identity to \(F^{-1}\circ F\) shows that \(\dd F_p\) is invertible,
with inverse \(\dd(F^{-1})_{F(p)}\).

\subsection{Changing coordinates at a fixed point}

Now consider two charts on the \emph{same} manifold, with coordinates
related by \(y^a=y^a(x)\). The point \(p\) and the vector \(v\in T_pM\)
remain fixed. Only their coordinate descriptions change. The chain rule
for partial derivatives gives
\begin{equation}
 \left.\frac{\partial}{\partial y^a}\right|_p
 =\frac{\partial x^i}{\partial y^a}(p)\,\partial_i|_p,
 \qquad
 v'^a=\frac{\partial y^a}{\partial x^i}(p)\,v^i.
 \label{eq:vector-coordinate-change}
\end{equation}
The second formula follows by writing \(v\) in both bases, or by applying
\(v\) to \(y^a\). The basis and its component column transform inversely,
as in SR1. A Lorentz coordinate change has a constant Jacobian and gives
the familiar law \(v'^\mu=\Lambda^\mu{}_\nu v^\nu\).

A coordinate change here is not a map moving the event \(p\) to a new
event. Its Jacobian converts between bases of one tangent space. By
contrast, the differential of a general map \(F\) sends a vector at
\(p\) to a vector at \(F(p)\).

The first derivative of a curve transforms by the vector law. A second
differentiation also differentiates the Jacobian. Along the curve,
\begin{equation}
 \ddot y^a
 =\frac{\partial y^a}{\partial x^i}\ddot x^i
  +\frac{\partial^2y^a}{\partial x^i\partial x^j}
       \dot x^i\dot x^j.
 \label{eq:coordinate-acceleration}
\end{equation}
All derivatives of the coordinate transition are evaluated at the current
point on the curve. The second term prevents \(\ddot x^i\) from being
the components of a vector under general coordinate changes. A connection
will provide the correction needed to define covariant acceleration.

\begin{studentwork}[3.6cm]{second derivatives under a coordinate change}
Derive \eqref{eq:coordinate-acceleration} and apply it to
\(x^1(t)=t\), \(x^2(t)=0\), with
\(y^1=x^1\), \(y^2=x^2+(x^1)^2\).
Then derive the transformation law of \(\partial_i\partial_j f\)
for a scalar field and determine why it is not generally a
covariant two-tensor.
\end{studentwork}

\section{Cotangent spaces, bundles, and fields}
\label{sec:bundles}
\slideref{10}

\subsection{Covectors and the differential of a function}

The tangent space \(T_pM\) is a vector space, so its dual is defined just
as in SR1. The \emph{cotangent space} is
\[
 T_p^*M=(T_pM)^*=\operatorname{Hom}_{\R}(T_pM,\R).
\]
Its elements are covectors at \(p\). A covector takes a tangent vector as
its input and returns a real number. This differs from a tangent vector
viewed as a derivation, whose input is a smooth function germ.

\begin{definition}{Differential of a scalar field}{scalar-differential}
For a smooth function \(f\) defined near \(p\), its \emph{differential at
\(p\)} is the covector
\[
 (\dd f)_p(v)=v(f),\qquad v\in T_pM.
\]
For coordinate functions \(x^1,\ldots,x^n\), their differentials satisfy
\begin{equation}
 \dd x^i|_p(\partial_j|_p)=\delta^i{}_j,
 \qquad
 (\dd f)_p=\partial_i f(p)\,\dd x^i|_p.
 \label{eq:df}
\end{equation}
Thus \(\dd x^1|_p,\ldots,\dd x^n|_p\) form the dual coordinate basis.
\end{definition}

The first identity follows by evaluating \(\partial_j|_p\) on \(x^i\).
For the second, both sides evaluated on \(v=v^j\partial_j|_p\) give
\(v^j\partial_j f(p)\), so they are the same covector. In particular,
\(\dd x^i|_p\) extracts the component \(v^i\) of a tangent vector.
For a curve velocity,
\[
 (\dd f)_p(\dot c(t_0))
 =\left.\frac{\dd}{\dd t}\right|_{t=t_0}f(c(t)),
 \qquad c(t_0)=p.
\]
The differential records all directional rates of change of \(f\) at the
point. It is not tied to one particular curve.

This agrees with the differential of the map \(f:M\to\R\) from the
previous section. That differential takes values in \(T_{f(p)}\R\),
which is canonically identified with \(\R\) by its usual coordinate.
Under this identification it is precisely the covector \((\dd f)_p\).

\slideref{11}
For \(v=v^i\partial_i|_p\) and \(\omega=\omega_i\dd x^i|_p\),
evaluation gives \(\omega(v)=\omega_i v^i\). Under \(y^a=y^a(x)\),
\begin{equation}
 \dd y^a=\frac{\partial y^a}{\partial x^i}\dd x^i,
 \qquad
 \omega'_a=\frac{\partial x^i}{\partial y^a}\omega_i.
 \label{eq:covector-coordinate-change}
\end{equation}
The first relation follows from \eqref{eq:df}. Expressing \(\omega\) in
both dual bases gives the second. Together with the vector transformation
law, it gives \(\omega'_a v'^a=\omega_i v^i\).

Evaluation also gives the canonical identification
\[
 T_pM\longrightarrow T_p^{**}M,
 \qquad v\longmapsto\bigl[\omega\longmapsto\omega(v)\bigr].
\]
This identifies a vector with a linear functional on \emph{covectors}.
It does not identify \(T_pM\) with \(T_p^*M\). A metric will later
provide that additional identification. In particular, \(\dd f\) is a
covector field, whereas its gradient will be a metric-dependent vector
field.

\subsection{Assembling the tangent spaces}

To describe vectors at all points of \(M\), we must retain both the base
point and its tangent vector. The tangent bundle collects these data in
one space. Vectors based at different points remain distinct elements,
even when their coordinate component lists happen to agree.


\begin{definition}{Tangent and cotangent bundles}{tangent-bundle}
The \emph{tangent bundle} and \emph{cotangent bundle} are the disjoint
unions
\[
 TM=\bigsqcup_{p\in M}T_pM,
 \qquad T^*M=\bigsqcup_{p\in M}T_p^*M.
\]
An element \(v_p\in TM\) records a point \(p\) and a vector in \(T_pM\).
The \emph{bundle projection} sends this element to its base point,
\(\pi(v_p)=p\). The set \(\pi^{-1}(p)=T_pM\) is the \emph{fiber}
over \(p\). The cotangent projection and its fibers are defined in the
same way.
\end{definition}

A chart \((U,\varphi)\), with coordinates \(x^1,\ldots,x^n\), supplies
coordinates on \(\pi^{-1}(U)\subset TM\). Expand
\(v_p=v^i\partial_i|_p\) and assign the tuple
\[
 v_p\longmapsto
 (x^1(p),\ldots,x^n(p),v^1,\ldots,v^n).
\]
The first \(n\) numbers specify the point. The remaining \(n\) numbers
specify the vector in the coordinate basis at that point. For
\(\omega_p=\omega_i\dd x^i|_p\), the analogous cotangent coordinates
are \((x^1(p),\ldots,x^n(p),\omega_1,\ldots,\omega_n)\).

On an overlap, the induced coordinate changes are
\[
 \begin{aligned}
 (x^i,v^i)&\longmapsto
 \left(y^a(x),\frac{\partial y^a}{\partial x^i}(x)v^i\right),\\
 (x^i,\omega_i)&\longmapsto
 \left(y^a(x),\frac{\partial x^i}{\partial y^a}(y(x))\omega_i\right).
 \end{aligned}
\]
Both are smooth, and both are linear in the vector or covector components
when the base point is fixed. They equip \(TM\) and \(T^*M\) with smooth
manifold structures. Their dimensions are
\[
 \dim T_pM=\dim T_p^*M=n,
 \qquad \dim TM=\dim T^*M=2n.
\]
The dimension \(2n\) counts both position and fiber coordinates. An
individual fiber still has dimension \(n\).

Both bundles are locally products, and their changes of fiber coordinates
are linear. The following definition isolates this common structure.

\begin{definition}{Smooth vector bundle}{vector-bundle}
A smooth \emph{vector bundle of rank \(r\)} over \(M\) consists of a
smooth manifold \(E\), a smooth surjection \(\pi:E\to M\), and an
\(r\)-dimensional vector-space structure on each fiber
\(E_p=\pi^{-1}(p)\). Every \(p\in M\) has an open neighborhood \(U\)
and a diffeomorphism
\[
 \Phi:\pi^{-1}(U)\longrightarrow U\times\R^r
\]
that leaves the base point unchanged and restricts to a linear
isomorphism \(E_q\to\{q\}\times\R^r\) for every \(q\in U\).
The map \(\Phi\) is a \emph{local trivialization}.
\end{definition}

A local trivialization represents each fiber over \(U\) by a copy of
\(\R^r\), with the representation varying smoothly with the point.
The constructions above make \(TM\) and \(T^*M\) rank-\(n\) vector
bundles. The existence of these local product descriptions does not
assert that either bundle is a product over all of \(M\).

\subsection{Sections and tensor fields}

A field chooses one element from each fiber. It is described by a map
into the total space whose projection returns the point where the field
is evaluated.

\begin{definition}{Section, vector field, and one-form}{field}
A \emph{section} of \(\pi:E\to M\) is a map \(s:M\to E\) satisfying
\(\pi\circ s=\id_M\). The space of smooth sections is denoted by
\(\Gamma(E)\).

A \emph{smooth vector field} is a smooth section of \(TM\), and a
\emph{smooth one-form} is a smooth section of \(T^*M\). Their spaces are
\[
 \X(M)=\Gamma(TM),\qquad \Omega^1(M)=\Gamma(T^*M).
\]
\end{definition}

Locally, a vector field and a one-form have the expressions
\(X=X^i\partial_i\) and \(\alpha=\alpha_i\dd x^i\). Smoothness means
that their coefficient functions are smooth. A tangent vector \(X_p\)
differentiates a function to a number. A vector field performs this
operation at every point, producing another smooth function,
\[
 X(f)(p)=X_p(f),\qquad X(f)=X^i\partial_i f.
\]
It obeys \(X(fh)=X(f)h+fX(h)\).

Fields can be added and multiplied by real constants. They can also be
multiplied by smooth functions, with \((fX)_p=f(p)X_p\). Thus \(\X(M)\)
is a real vector space and a \emph{module over \(\Cinf(M)\)}. Here
``module'' means that the multipliers may be smooth functions rather
than only real numbers. Linearity with respect to such multipliers is
stronger than real linearity. This distinction will be essential when
defining derivatives of fields.

A \emph{local frame} consists of smooth vector fields
\(e_1,\ldots,e_n\) that form a basis at every point of an open set.
Their pointwise dual one-forms form the \emph{dual coframe}.
Coordinate fields give one example of a local frame. An arbitrary smooth
frame need not arise from coordinates.

A diffeomorphism \(F:M\to N\) pushes a vector field forward by
\[
 (F_*X)_{F(p)}=\dd F_p(X_p).
\]
Each target point has exactly one preimage, so this defines a smooth
field on \(N\). For a general smooth map, different preimages may
produce different vectors at the same target point. The pointwise
differential therefore does not automatically define a vector field
on the target.

More generally, a tensor field of type \((r,s)\) is a smooth section of
\[
 T^r_sM=(TM)^{\otimes r}\otimes(T^*M)^{\otimes s}.
\]
The fiber over \(p\) is
\((T_pM)^{\otimes r}\otimes(T_p^*M)^{\otimes s}\). All its factors
refer to the same point. Its components obey the tensor transformation
law from SR1, with the coordinate Jacobian replacing the constant
Lorentz matrix.

For example, a smooth family of linear maps \(A_p:T_pM\to T_pM\) is
a tensor field of type \((1,1)\), written
\[
 A=A^i{}_j\,\partial_i\otimes\dd x^j,
 \qquad A(X)=A^i{}_jX^j\partial_i.
\]
Under \(y^a=y^a(x)\), its components at the same point satisfy
\[
 A'^a{}_b(p)=
 \frac{\partial y^a}{\partial x^i}(p)
 \frac{\partial x^j}{\partial y^b}(p)A^i{}_j(p).
\]
Here \(A^i{}_j\) and \(A'^a{}_b\) denote the component functions in
the two charts. The identity map has components \(\delta^i{}_j\) in
every chart.

A useful test for tensoriality is linearity over \(\Cinf(M)\) in
each field argument. Expand the fields in a local basis. Such linearity
allows their component functions to be taken outside the operation,
without differentiating them. For example, if a smooth operation
\(B(X,Y)\) has this property in both arguments, then locally
\[
 B(X,Y)_p=X^i(p)Y^j(p)\,B(\partial_i,\partial_j)_p.
\]
Its value depends only on \(X_p\) and \(Y_p\), not on derivatives
of the fields. A scalar-valued operation of this kind is a covariant
two-tensor. A vector-valued one is a tensor of type \((1,2)\).

\section{Tangent geometry of the two-sphere}
\label{sec:sphere-tangent}
\slideref{12}

The sphere lets us compare the intrinsic definitions with familiar vectors
in three-dimensional space. Its tangent vectors can be represented by
ambient arrows, provided their base points are retained.

\subsection{Intrinsic vectors and the ambient tangent plane}

Let \(\iota:S^2\hookrightarrow\R^3\) be the inclusion. A curve on the
sphere is also a curve in \(\R^3\), and \(\dd\iota_p\) sends its
intrinsic velocity to its ambient velocity. Using Cartesian coordinates
to identify \(T_{\iota(p)}\R^3\) with \(\R^3\), we obtain the usual
tangent plane as a vector subspace.

\begin{proposition}{Tangent space of the unit sphere}{sphere-tangent}
For every \(p\in S^2\),
\[
 \dd\iota_p(T_pS^2)=p^\perp
 =\{v\in\R^3:p\cdot v=0\}.
\]
With this identification, the tangent bundle can be written as
\[
 TS^2=\{(p,v)\in\R^3\times\R^3:p\cdot p=1,\ p\cdot v=0\}.
\]
\end{proposition}

\begin{proof}
If \(c(0)=p\) is a curve on the sphere, differentiating
\(c(t)\cdot c(t)=1\) gives \(p\cdot\dot c(0)=0\). Conversely, for
\(v\in p^\perp\), the curve
\[
 c(t)=\frac{p+tv}{\sqrt{1+t^2\lVert v\rVert^2}}
\]
lies on the sphere and has \(c(0)=p\), \(\dot c(0)=v\). This also covers
\(v=0\). The image of \(\dd\iota_p\) is therefore the two-dimensional
space \(p^\perp\). Since \(\dim T_pS^2=2\), the differential is an
isomorphism onto this image.
\end{proof}

The affine tangent plane drawn through the point is \(p+p^\perp\). The
vector space representing tangent vectors is \(p^\perp\). Keeping the base
point in the pair \((p,v)\) distinguishes the fibers even when their ambient
vector representations happen to coincide.

\subsection{Spherical coordinates and coordinate vectors}

On the sphere away from the poles and a chosen longitude cut, use
\[
 \sigma(\vartheta,\varphi)=
 (\sin\vartheta\cos\varphi,\sin\vartheta\sin\varphi,\cos\vartheta),
\]
where \(0<\vartheta<\pi\) and
\(\varphi_0<\varphi<\varphi_0+2\pi\). The inverse of this parametrization
is the chart. The ambient representatives of its coordinate vectors are
\begin{align}
 e_\vartheta&=\frac{\partial\sigma}{\partial\vartheta}
 =(\cos\vartheta\cos\varphi,\cos\vartheta\sin\varphi,-\sin\vartheta),
 \label{eq:sphere-etheta}\\
 e_\varphi&=\frac{\partial\sigma}{\partial\varphi}
 =\sin\vartheta(-\sin\varphi,\cos\varphi,0).
 \label{eq:sphere-ephi}
\end{align}
These represent \(\partial_\vartheta\) and \(\partial_\varphi\), respectively.
The coordinate vectors are tangent and linearly independent on the stated chart. Hence
\[
 V=V^\vartheta\partial_\vartheta+V^\varphi\partial_\varphi
\]
has ambient representative
\(V^\vartheta e_\vartheta+V^\varphi e_\varphi\). The tuple
\((\vartheta,\varphi,V^\vartheta,V^\varphi)\) gives local coordinates on
\(TS^2\).

At either pole, the formal expression for \(e_\varphi\) vanishes. The tangent
space has not lost a dimension. Longitude is not a coordinate at a pole, so
this particular chart and its frame do not apply there.

\begin{studentwork}[3.7cm]{a global vector field and a local coordinate frame}
For \(K(x,y,z)=(-y,x,0)\), prove tangency to \(S^2\), find its angular
coordinate components, and locate its zeros. Explain why \(K\) is smooth
at the poles although the angular coordinate frame is not defined there.
\end{studentwork}

\subsection{Comparing vectors at different points}
\slideref{13}

Changing the basis at \(p\) keeps the vector in the same space
\(T_pS^2\). Comparing a vector at \(p\) with one at another point \(q\)
is a different operation. The second vector belongs to \(T_qS^2\), so
\(V_q-V_p\) is not defined by the manifold structure alone.

The embedding lets us subtract the ambient representatives in \(\R^3\).
But differentiating a tangent field in \(\R^3\) need not give a tangent
vector. Its derivative may contain a component normal to the sphere.
To obtain a derivative within the tangent spaces, a rule for comparing
the fibers along a curve must be specified. This is the role of a
connection.


\section{Connections and parallel transport}
\label{sec:connections}
\slideref{14}

\subsection{Covariant differentiation}

For a scalar field, nearby values can be subtracted because they all
belong to \(\R\). For a vector field \(Y\), the values \(Y_p\) and
\(Y_q\) belong to different tangent spaces. A derivative must specify
how these values are compared. An affine connection supplies such a
rule.

Differentiating the component functions alone does not solve the problem.
If \(Y'^a=(\partial y^a/\partial x^j)Y^j\), then
\[
 \partial'_bY'^a
 =\frac{\partial x^i}{\partial y^b}
   \frac{\partial y^a}{\partial x^j}\partial_iY^j
  +\frac{\partial x^i}{\partial y^b}
   \frac{\partial^2y^a}{\partial x^i\partial x^j}Y^j.
\]
The second term comes from differentiating the change-of-basis matrix.
It is absent from a tensor transformation law. A connection adds terms
that compensate for this dependence on the chosen frame. Choosing the
connection is additional geometric data, not just a change of notation.

\begin{definition}{Affine connection}{connection}
An \emph{affine connection} on \(TM\) is a map
\[
 \nabla:\X(M)\times\X(M)\longrightarrow\X(M),
 \qquad (X,Y)\longmapsto\nabla_XY,
\]
satisfying
\[
 \begin{aligned}
 \nabla_{fX+hY}Z&=f\nabla_XZ+h\nabla_YZ,\\
 \nabla_X(aY+bZ)&=a\nabla_XY+b\nabla_XZ,\\
 \nabla_X(fY)&=X(f)Y+f\nabla_XY,
 \end{aligned}
\]
for smooth fields \(X,Y,Z\), functions \(f,h\in\Cinf(M)\), and constants
\(a,b\in\R\). The field \(\nabla_XY\) is the \emph{covariant derivative
of \(Y\) along \(X\)}.
\end{definition}

The first identity is linearity over \(\Cinf(M)\) in \(X\). At a point,
only the vector \(X_p\) is needed to specify the differentiation
direction. The second and third identities say that differentiation in
\(Y\) is real-linear and obeys a product rule. It may therefore depend
on the first derivatives of \(Y\), not just its value at the point.

To work in a chart, we need the derivative to depend only on the fields
near the point. This locality follows from the axioms. Suppose \(Y\)
vanishes near \(p\). Choose a smooth cutoff \(\chi\) equal to one near
\(p\), with support inside that neighborhood. It can be constructed in
a coordinate ball and extended by zero. Since \(\chi Y=0\), the product
rule at \(p\) gives \((\nabla_XY)_p=0\). First-slot linearity gives
the same conclusion when \(X\) vanishes near \(p\). Thus local fields
may be extended using cutoffs, and the derivative near \(p\) is
independent of the extension. We can therefore apply \(\nabla\) to
local coordinate fields.

Define the \emph{connection coefficients} by
\begin{equation}
 \nabla_{\partial_i}\partial_j=\Gamma^k{}_{ij}\partial_k.
 \label{eq:gamma-convention}
\end{equation}
The first lower index is the differentiation direction. The second
identifies the basis vector being differentiated. For
\(X=X^i\partial_i\) and \(Y=Y^j\partial_j\), the axioms give
\[
 \begin{aligned}
 \nabla_XY
 &=X^i\nabla_{\partial_i}(Y^j\partial_j)\\
 &=X^i(\partial_iY^j)\partial_j
   +X^iY^j\Gamma^k{}_{ij}\partial_k.
 \end{aligned}
\]
Renaming the summed indices yields
\begin{equation}
 \nabla_XY
 =X^i\bigl(\partial_iY^k+\Gamma^k{}_{ij}Y^j\bigr)\partial_k.
 \label{eq:covariant-vector}
\end{equation}
The first term differentiates the components. The second accounts for
the covariant derivatives of the basis vectors. Their sum is a vector,
although the two terms do not separately transform as vectors.

For fixed \(Y\), the map \(X\mapsto\nabla_XY\) is linear over
\(\Cinf(M)\). It therefore defines a tensor field of type \((1,1)\),
\[
 \nabla Y=(\partial_iY^k+\Gamma^k{}_{ij}Y^j)
           \partial_k\otimes\dd x^i.
\]
The notation \(\nabla_iY^k\) denotes the component in parentheses. It
is not the ordinary derivative of the single scalar function \(Y^k\).

\begin{proposition}{Transformation of connection coefficients}{gamma-transformation}
For the coordinate change \(y^a=y^a(x)\), the coefficients of the same
connection satisfy
\begin{equation}
 \widetilde\Gamma^a{}_{bc}
 =\frac{\partial y^a}{\partial x^k}
  \frac{\partial x^i}{\partial y^b}
  \frac{\partial x^j}{\partial y^c}\Gamma^k{}_{ij}
 +\frac{\partial y^a}{\partial x^k}
  \frac{\partial^2x^k}{\partial y^b\partial y^c}.
 \label{eq:gamma-transform}
\end{equation}
The second term shows that connection coefficients are not the
components of a tensor.
\end{proposition}

To obtain this law, substitute
\(\partial'_c=(\partial x^j/\partial y^c)\partial_j\) into
\(\nabla_{\partial'_b}\partial'_c\). Differentiating the coefficient
\(\partial x^j/\partial y^c\) produces the second-derivative term.
Expressing the result in the primed basis gives
\eqref{eq:gamma-transform}. This extra term cancels the unwanted term
from the transformation of \(\partial_iY^k\).

\subsection{Differentiation along a curve}

Let \(\gamma:[a,b]\to M\) be a smooth curve. A \emph{vector field along
\(\gamma\)} assigns a vector \(V(t)\in T_{\gamma(t)}M\), with smooth
components in local coordinates. Write
\(V(t)=V^k(t)\partial_k|_{\gamma(t)}\). This need not come from one field
on \(M\). For example, a curve can pass through the same point at two
parameter values and carry different vectors on the two visits.

The same product rule used for \(\nabla_XY\) now differentiates
\(V^k(t)\), while the basis varies along \(\gamma\). The direction
components \(X^i\) are replaced by the curve velocity \(\dot\gamma^i\).

\begin{definition}{Covariant derivative along a curve}{curve-derivative}
The connection defines
\begin{equation}
 \frac{DV}{\dd t}
 =\left(\frac{\dd V^k}{\dd t}
        +\Gamma^k{}_{ij}(\gamma(t))\dot\gamma^i V^j\right)
       \partial_k|_{\gamma(t)}.
 \label{eq:curve-covariant}
\end{equation}
The field \(V\) is \emph{parallel along \(\gamma\)} when
\(DV/\dd t=0\).
\end{definition}

The transformation law of the connection coefficients makes this
expression coordinate independent. If \(V(t)=Y_{\gamma(t)}\) for a
local vector field \(Y\), it agrees with \(\nabla_{\dot\gamma}Y\).
Taking \(V=\dot\gamma\) gives the covariant acceleration,
\[
 \frac{D\dot\gamma}{\dd t}
 =\bigl(\ddot\gamma^k+
        \Gamma^k{}_{ij}\dot\gamma^i\dot\gamma^j\bigr)
      \partial_k|_{\gamma(t)}.
\]
The connection term corrects the nonvectorial transformation of ordinary
coordinate acceleration found in Section~\ref{sec:differentials}.

\begin{theorem}{Parallel transport along a curve}{parallel-transport}
For each \(v\in T_{\gamma(a)}M\), there is a unique parallel field
along \(\gamma\) satisfying \(V(a)=v\). Its endpoint value defines a
linear isomorphism
\[
 P_\gamma:T_{\gamma(a)}M\longrightarrow T_{\gamma(b)}M,
 \qquad P_\gamma(v)=V(b),
\]
called \emph{parallel transport along \(\gamma\)}.
\end{theorem}

\begin{proof}
In a chart, the parallel condition is the linear system of ordinary
differential equations
\[
 \dot V^k=-\Gamma^k{}_{ij}(\gamma(t))\dot\gamma^iV^j.
\]
Existence and uniqueness hold on each compact subinterval contained in
one chart. A finite subdivision of \([a,b]\) allows these solutions
to be continued along the whole curve and joined uniquely on overlaps.
Linearity of the equations makes the endpoint map linear. For the
reversed curve \(\bar\gamma(t)=\gamma(a+b-t)\), the evolution is
reversed, giving \(P_{\bar\gamma}=P_\gamma^{-1}\).
\end{proof}

Parallel transport now provides a comparison of vectors in different
tangent spaces. To compare \(v\in T_{\gamma(a)}M\) with
\(w\in T_{\gamma(b)}M\), first transport \(v\). The difference
\(w-P_\gamma v\) is defined because both terms lie in the endpoint
tangent space. The result can depend on the chosen path, not only on
its endpoints.

The same comparison explains covariant differentiation. Fix \(t_0\)
and let \(P_{t_0+h\to t_0}\) denote transport back along \(\gamma\).
For a field \(V\) along the curve,
\[
 \left.\frac{DV}{\dd t}\right|_{t_0}
 =\lim_{h\to0}
 \frac{P_{t_0+h\to t_0}V(t_0+h)-V(t_0)}{h}.
\]
All vectors in this difference quotient belong to \(T_{\gamma(t_0)}M\).
Expanding the transport equation to first order in \(h\) recovers
\eqref{eq:curve-covariant}.

\begin{examplebox}{Tangential differentiation on the sphere}{sphere-connection}
For \(\gamma(t)\in S^2\), represent a tangent field by an ambient
vector satisfying \(\gamma\cdot V=0\). Using the Euclidean scalar
product in \(\R^3\), project its derivative onto the tangent plane,
\[
 \frac{DV}{\dd t}
 =\dot V-(\gamma\cdot\dot V)\gamma
 =\dot V+(\dot\gamma\cdot V)\gamma.
\]
The second equality follows by differentiating
\(\gamma\cdot V=0\). This rule defines a connection on the sphere.
Along the equator \(\gamma(t)=(\cos t,\sin t,0)\), the field
\(V=\dot\gamma\) has \(\dot V=-\gamma\) and \(DV/\dd t=0\).
Its ambient derivative is nonzero, but entirely normal to the sphere.
\end{examplebox}

\begin{examplebox}{The Euclidean connection in polar coordinates}{polar-connection}
Choose the connection on the plane for which the Cartesian basis fields
are parallel. On a polar chart,
\[
 \partial_r=\cos\varphi\,\partial_x+\sin\varphi\,\partial_y,
 \qquad
 \partial_\varphi=-r\sin\varphi\,\partial_x+r\cos\varphi\,\partial_y.
\]
The product rule gives
\[
 \nabla_{\partial_r}\partial_r=0,\qquad
 \nabla_{\partial_r}\partial_\varphi
 =\nabla_{\partial_\varphi}\partial_r=\frac1r\partial_\varphi,
 \qquad
 \nabla_{\partial_\varphi}\partial_\varphi=-r\partial_r.
\]
Thus the nonzero coefficients are
\[
 \Gamma^r{}_{\varphi\varphi}=-r,\qquad
 \Gamma^\varphi{}_{r\varphi}=\Gamma^\varphi{}_{\varphi r}=\frac1r.
\]
The Cartesian coefficients vanish and the polar coefficients do not,
but both describe the same connection.
\end{examplebox}

\begin{studentwork}[3.8cm]{parallelism and varying components}
Verify the polar-coordinate derivatives above. Express \(\partial_x\)
in the polar frame and show that
\(\nabla_{\partial_r}\partial_x=\nabla_{\partial_\varphi}\partial_x=0\).
Explain why a parallel field need not have constant coordinate components.
\end{studentwork}

No metric was required in the definition of an affine connection.
A metric selects a particular connection once metric compatibility and
zero torsion are imposed. Appendix~\ref{sec:levi-civita} develops these
conditions after the metric and Lie bracket have been introduced.

\section{Differential forms and exterior algebra}
\label{sec:forms}
\slideref{15}

A one-form acts linearly on a tangent vector. A differential form of
higher degree acts multilinearly on several tangent vectors and is
\emph{alternating}, meaning that interchanging two arguments reverses
its sign. All the arguments are vectors at the same point.
This antisymmetry will allow a derivative in which the extra terms
arising from coordinate changes cancel.

\begin{definition}{Differential form}{form}
For \(k\geq1\), the space \(\Lambda^kT_p^*M\) consists of alternating
multilinear maps
\[
 \alpha_p:(T_pM)^k\longrightarrow\R.
\]
A \emph{smooth differential \(k\)-form} assigns such a map to each
\(p\in M\), with smooth components in local coordinates. Equivalently,
it is a smooth section of \(\Lambda^kT^*M\). We write
\[
 \Omega^k(M)=\Gamma(\Lambda^kT^*M),\qquad
 \Omega^0(M)=\Cinf(M).
\]
Thus a zero-form is a smooth function.
\end{definition}

\subsection{The wedge product}

For one-forms \(\alpha\) and \(\beta\), their \emph{wedge product} is
the two-form defined by
\[
 (\alpha\wedge\beta)(v,w)
 =\alpha(v)\beta(w)-\alpha(w)\beta(v).
\]
The subtraction makes the result alternating. In particular,
\[
 \dd x^i\wedge\dd x^j=-\dd x^j\wedge\dd x^i,
 \qquad \dd x^i\wedge\dd x^i=0.
\]
More generally, the wedge product of \(k\) covectors at one point is
\[
 (\alpha^1\wedge\cdots\wedge\alpha^k)(v_1,\ldots,v_k)
 =\det\bigl(\alpha^i(v_j)\bigr).
\]
Such products span \(\Lambda^kT_p^*M\), but not every form is a single
product of one-forms.

For \(\alpha\in\Omega^k(M)\) and \(\beta\in\Omega^\ell(M)\),
evaluate the forms on separate groups of vectors, then antisymmetrize
over all \(k+\ell\) inputs. Write \(S_m\) for the permutations of
\(\{1,\ldots,m\}\), and \(\operatorname{sgn}(\sigma)\) for the sign
of a permutation \(\sigma\). The resulting product is
\[
 \begin{aligned}
 (\alpha\wedge\beta)(v_1,\ldots,v_{k+\ell})
 =\frac1{k!\,\ell!}\sum_{\sigma\in S_{k+\ell}}
 &\operatorname{sgn}(\sigma)\,
  \alpha(v_{\sigma(1)},\ldots,v_{\sigma(k)})\\
 &\times\beta(v_{\sigma(k+1)},\ldots,v_{\sigma(k+\ell)}).
 \end{aligned}
\]
Each form is already alternating in its own inputs. Permutations within
these two groups therefore repeat each contribution \(k!\,\ell!\)
times, which explains the normalization. The formula defines a
\((k+\ell)\)-form without first expressing the factors as sums of
products of one-forms. Wedge multiplication is associative and satisfies
\[
 \alpha\wedge\beta=(-1)^{k\ell}\beta\wedge\alpha.
\]
This is \emph{graded commutativity}. The sign comes from moving \(k\)
one-form factors past \(\ell\) others, requiring \(k\ell\) swaps.
Degree-zero factors act by ordinary multiplication of a form by a function.

\subsection{Components and normalization}

The components of a \(k\)-form are defined by evaluation on coordinate
basis vectors,
\[
 \alpha_{i_1\cdots i_k}
 =\alpha(\partial_{i_1},\ldots,\partial_{i_k}).
\]
They are completely antisymmetric. The products
\(\dd x^{i_1}\wedge\cdots\wedge\dd x^{i_k}\), with
\(i_1<\cdots<i_k\), form a basis of \(\Lambda^kT_p^*M\) at each point.
Consequently,
\begin{equation}
 \begin{aligned}
 \alpha
 &=\sum_{i_1<\cdots<i_k}\alpha_{i_1\cdots i_k}
       \dd x^{i_1}\wedge\cdots\wedge\dd x^{i_k}\\
 &=\frac1{k!}\alpha_{i_1\cdots i_k}
       \dd x^{i_1}\wedge\cdots\wedge\dd x^{i_k}.
 \end{aligned}
 \label{eq:form-components}
\end{equation}
The second line sums over all index values. Permuting distinct indices
changes the signs of both the component and the wedge product, so the
same contribution occurs \(k!\) times. This explains the factor
\(1/k!\). For example,
\[
 F=\frac12F_{ij}\dd x^i\wedge\dd x^j,
 \qquad F(\partial_i,\partial_j)=F_{ij}.
\]

There are \(\binom nk\) increasing index lists of length \(k\). Hence
\[
 \dim\Lambda^kT_p^*M=\binom nk,
 \qquad
 \dim\left(\bigoplus_{k=0}^n\Lambda^kT_p^*M\right)=2^n.
\]
These are dimensions of spaces of forms at \emph{one point}. The space
\(\Omega^k(M)\) of smooth form fields is generally infinite-dimensional
because its coefficients are functions. For \(k>n\), alternating
\(k\)-forms vanish: any \(k\) tangent vectors are linearly dependent.

No metric has been used. The determinant formula assigns a signed
number once the covectors have been specified. It does not by itself
supply the lengths or angles of the input vectors. Also, graded
commutativity forces \(\alpha\wedge\alpha=0\) for odd-degree forms,
but imposes no such condition for even degree.

\begin{studentwork}[3.2cm]{alternating products in four dimensions}
For \(F=\dd x^1\wedge\dd x^2+\dd x^3\wedge\dd x^4\) on \(\R^4\),
compute \(F\wedge F\). Deduce that \(F\) cannot be written as
\(\alpha\wedge\beta\) for two one-forms.
\end{studentwork}

\section{Exterior differentiation and pullback}
\label{sec:exterior}
\slideref{16}

\subsection{The exterior derivative}

For a scalar field, \(\dd f\) is a one-form that records directional
rates of change. The \emph{exterior derivative} extends this operation
to forms of every degree. It takes a \(k\)-form to a \((k+1)\)-form
and requires neither a metric nor a connection.

\begin{theorem}{Exterior derivative}{exterior-derivative}
On a smooth manifold there is a unique family of real-linear operators
\[
 \dd:\Omega^k(M)\longrightarrow\Omega^{k+1}(M)
\]
that agrees with the differential on functions and satisfies
\[
 \dd^2=0,\qquad
 \dd(\alpha\wedge\beta)
 =\dd\alpha\wedge\beta+(-1)^k\alpha\wedge\dd\beta,
 \qquad \alpha\in\Omega^k(M).
\]
In a chart, let \(I=(i_1,\ldots,i_k)\) range over increasing index
lists and write
\(\dd x^I=\dd x^{i_1}\wedge\cdots\wedge\dd x^{i_k}\).
For a coordinate expansion
\(\alpha=\sum_I\alpha_I\,\dd x^I\), its local formula is
\begin{equation}
 \dd\alpha
 =\sum_I\dd\alpha_I\wedge\dd x^I
 =\sum_{I,j}\partial_j\alpha_I\,\dd x^j\wedge\dd x^I.
 \label{eq:exterior-local}
\end{equation}
\end{theorem}

The local rule differentiates each coefficient and places its
differential before the existing wedge factors. The coordinate one-forms
contribute no extra terms because \(\dd(\dd x^i)=\dd^2x^i=0\).
For a one-form,
\[
 \dd(\alpha_i\dd x^i)
 =\partial_j\alpha_i\,\dd x^j\wedge\dd x^i.
\]
Renaming the summed indices and using antisymmetry gives
\begin{equation}
 \dd\alpha
 =\frac12(\partial_i\alpha_j-\partial_j\alpha_i)
           \dd x^i\wedge\dd x^j.
 \label{eq:d-oneform}
\end{equation}
Thus only the antisymmetric part of the coefficient derivatives remains.
For \(F=\tfrac12F_{ij}\dd x^i\wedge\dd x^j\), the corresponding
components are
\[
 (\dd F)_{ijk}
 =\partial_iF_{jk}+\partial_jF_{ki}+\partial_kF_{ij}.
\]

\begin{examplebox}{Exterior derivative of a one-form}{exterior-example}
On \(\R^2\), let \(\alpha=x\,\dd y-y\,\dd x\). Then
\[
 \dd\alpha=\dd x\wedge\dd y-\dd y\wedge\dd x
           =2\,\dd x\wedge\dd y.
\]
For tangent vectors \(v,w\),
\((\dd\alpha)(v,w)=2(v^xw^y-v^yw^x)\). The result is an
alternating bilinear form.
\end{examplebox}

To construct \(\dd\), use \eqref{eq:exterior-local} in each chart.
The chain rule gives \(\dd y^a=(\partial_i y^a)\dd x^i\) on an
overlap. Applying the local rule once more gives
\[
 \dd(\dd y^a)
 =\partial_j\partial_i y^a\,\dd x^j\wedge\dd x^i=0,
\]
since the coefficient is symmetric in \(i,j\) and the wedge product
is antisymmetric. The local operators therefore agree on functions and
on coordinate differentials. The product rule makes them agree on every
form. The same cancellation of mixed partial derivatives gives
\(\dd^2=0\).

For uniqueness, the product rule implies locality by a cutoff argument,
as for connections. In a chart, the operator is fixed on functions and
has \(\dd(\dd x^i)=0\). Applying its product rule to the coordinate
expansion of any form then forces \eqref{eq:exterior-local}.

\begin{studentwork}[3.3cm]{the identity \(\dd^2=0\)}
Verify \(\dd^2f=0\) directly from the coordinate formula for a smooth
function. Extend the calculation to an arbitrary differential form.
\end{studentwork}

A form is \emph{closed} when \(\dd\alpha=0\). A form of positive
degree is \emph{exact} when \(\alpha=\dd\beta\) for a globally defined
form \(\beta\). Every exact form is closed because \(\dd^2=0\).
The converse need not hold globally. An \(n\)-form on an
\(n\)-manifold is always closed because there are no nonzero
\((n+1)\)-forms. This does not imply that it is exact.

\subsection{Pullback of forms}

A smooth map \(F:M\to N\) sends tangent vectors forward by its
differential. A form on \(N\) can then evaluate these image vectors.
This defines a form on \(M\), called its \emph{pullback}. The form
moves in the opposite direction to the map, but no inverse of \(F\)
is required.

\begin{definition}{Pullback}{pullback}
For \(F:M\to N\) smooth and \(\alpha\in\Omega^k(N)\), define
\(F^*\alpha\in\Omega^k(M)\) by
\[
 (F^*\alpha)_p(v_1,\ldots,v_k)
 =\alpha_{F(p)}(\dd F_pv_1,\ldots,\dd F_pv_k).
\]
For a function, the definition is \(F^*f=f\circ F\).
\end{definition}

For a one-form, first send \(v\in T_pM\) to \(\dd F_pv\), then
evaluate \(\alpha_{F(p)}\) on that vector. In target coordinates \(y^a\), it gives the practical
formulas
\[
 F^*(\dd y^a)=\dd(y^a\circ F),\qquad
 F^*(\alpha_a\dd y^a)=(\alpha_a\circ F)\,\dd(y^a\circ F).
\]
Both the coefficients and the coordinate differentials must be pulled
back. The active scalar-field transformation in SR1,
\(f\mapsto f\circ F^{-1}\), is therefore \((F^{-1})^*f\) when
\(F\) is a diffeomorphism. It should not be confused with \(F^*f\).

The same evaluation rule defines pullback for any covariant tensor.
For a covariant two-tensor \(h\),
\[
 (F^*h)_p(v,w)=h_{F(p)}(\dd F_pv,\dd F_pw).
\]
Even if \(h\) is nondegenerate, its pullback need not be. For example,
a nonzero vector in the kernel of \(\dd F_p\) is paired to zero with
every vector by \(F^*h\).

\begin{proposition}{Pullback and exterior differentiation}{naturality}
For smooth maps and differential forms,
\[
 F^*(\alpha\wedge\beta)=F^*\alpha\wedge F^*\beta,
 \qquad F^*(\dd\alpha)=\dd(F^*\alpha).
\]
If \(F:M\to N\) and \(G:N\to P\), then
\((G\circ F)^*=F^*\circ G^*\).
\end{proposition}

The wedge identity follows from the alternating evaluation rule, and
the composition identity follows from the chain rule for differentials.
For a function \(f\), the chain rule gives
\(F^*(\dd f)=\dd(f\circ F)\). Every form is locally a sum of terms
\(f\,\dd y^{i_1}\wedge\cdots\wedge\dd y^{i_k}\). Applying the
product rules extends the identity to all forms. Thus pulling back
and exterior differentiating can be performed in either order.

\begin{examplebox}{A smooth one-form on the sphere}{sphere-oneform}
Let \(\iota:S^2\hookrightarrow\R^3\) be the inclusion and define
\[
 \beta=\iota^*(x\,\dd y-y\,\dd x).
\]
Substituting the spherical parametrization gives
\[
 \beta=\sin^2\vartheta\,\dd\varphi,\qquad
 \dd\beta=2\sin\vartheta\cos\vartheta\,
                \dd\vartheta\wedge\dd\varphi.
\]
The pullback definition gives a smooth one-form on the entire sphere.
The angular expressions are valid only on their chart. In particular,
\(\dd\varphi\) itself is not defined at a pole, but this does not
prevent \(\beta\) from being smooth there.
\end{examplebox}

\begin{studentwork}[4.0cm]{pullback to a surface}
For \(F(u,v)=(u,v,u^2+v^2)\) and
\(\alpha=x\,\dd y+y\,\dd z+z\,\dd x\) on \(\R^3\), compute
\(F^*\alpha\) and verify \(\dd(F^*\alpha)=F^*(\dd\alpha)\) directly.
\end{studentwork}

\subsection{Local and global potentials}

A potential for a one-form \(\alpha\) is a smooth function \(f\) with
\(\alpha=\dd f\). Integration tests whether such a potential can be
defined globally. For an oriented smooth parametrized curve
\(\gamma:[a,b]\to M\),
\[
 \int_\gamma\alpha
 :=\int_a^b\alpha_{\gamma(t)}(\dot\gamma(t))\,\dd t.
\]
This is the integral of the pullback \(\gamma^*\alpha\) on the parameter
interval. A change of parameter preserving orientation leaves it
unchanged. For \(\alpha=\dd f\), the chain rule gives
\[
 \int_\gamma\dd f=f(\gamma(b))-f(\gamma(a)).
\]
In particular, an exact one-form integrates to zero around a closed curve.

On \(\R^2\setminus\{0\}\), consider
\[
 \alpha=\frac{-y\,\dd x+x\,\dd y}{x^2+y^2}.
\]
A coordinate calculation gives \(\dd\alpha=0\). In an angular chart,
\(\alpha=\dd\varphi\), so it has a local potential. But along the unit
circle \(\gamma(t)=(\cos t,\sin t)\), \(0\leq t\leq2\pi\),
\[
 \gamma^*\alpha=\dd t,\qquad \int_\gamma\alpha=2\pi.
\]
It follows that no globally defined real-valued function \(f\) satisfies
\(\alpha=\dd f\). The angle can be chosen smoothly on individual
angular charts, but not as a single real-valued angle on the punctured
plane. Thus local potentials do not always combine into a global one.

Exterior differentiation acts on alternating covariant tensors.
The next construction instead uses the local flow of a vector field
to differentiate vector fields and more general tensor fields, again
without choosing a connection.

\section{Vector fields, flows, and Lie derivatives}
\label{sec:lie}
\slideref{17}

\subsection{The Lie bracket}

A vector field \(X\) acts on smooth functions by \(f\mapsto X(f)\).
Given two fields \(X,Y\), we can differentiate a function first with
\(Y\) and then with \(X\), obtaining \(X(Y(f))\). Reversing the order
need not give the same result. Their difference defines the Lie bracket.

\begin{definition}{Lie bracket}{lie-bracket}
For smooth vector fields \(X,Y\), the \emph{Lie bracket} is the
vector field determined by
\[
 [X,Y](f)=X(Y(f))-Y(X(f)),\qquad f\in\Cinf(M).
\]
\end{definition}

The definition must produce a first-order derivation, not a general
second-order differential operator. To check this, expand in coordinates,
\[
 \begin{aligned}
 X(Y(f))-Y(X(f))
 &=(X^i\partial_iY^k-Y^i\partial_iX^k)\partial_kf\\
 &\quad +(X^iY^j-Y^iX^j)\partial_i\partial_jf.
 \end{aligned}
\]
The coefficient of the second derivatives is antisymmetric in \(i,j\),
while \(\partial_i\partial_jf\) is symmetric. Their contraction
vanishes, leaving
\begin{equation}
 [X,Y]=(X^i\partial_iY^k-Y^i\partial_iX^k)\partial_k.
 \label{eq:bracket-components}
\end{equation}
The resulting coefficients are smooth. Since the original definition
uses the actions of the fields on functions, the vector field does not
depend on a coordinate choice.

\begin{proposition}{The Lie algebra of vector fields}{lie-algebra}
The Lie bracket is real-bilinear and antisymmetric, and satisfies the
Jacobi identity
\[
 [X,[Y,Z]]+[Y,[Z,X]]+[Z,[X,Y]]=0.
\]
Thus \(\X(M)\), with this bracket, is a real Lie algebra.
\end{proposition}

A Lie algebra is a vector space with a bilinear bracket satisfying
antisymmetry and the Jacobi identity. Here that vector space is
\(\X(M)\), the space of vector fields, not the set of points of \(M\).
The bracket also satisfies
\begin{equation}
 [X,fY]=X(f)Y+f[X,Y],\qquad
 [fX,Y]=f[X,Y]-Y(f)X.
 \label{eq:bracket-function}
\end{equation}
It is therefore not linear over \(\Cinf(M)\) in either field argument.
The value \([X,Y]_p\) depends on how the fields vary near \(p\), not
just on \(X_p\) and \(Y_p\). Coordinate fields commute,
\([\partial_i,\partial_j]=0\), but general frame fields need not.

\begin{studentwork}[3.4cm]{commutators and the Jacobi identity}
Prove the Jacobi identity and \eqref{eq:bracket-function}. Explain why
the Lie bracket is not a tensor of type \((1,2)\).
\end{studentwork}

\subsection{Local flows}

A vector field specifies a tangent vector at each point. An integral
curve follows these vectors, with its velocity equal to the field at
its current position. In coordinates, this means solving
\[
 \frac{\dd x^i(\gamma(t))}{\dd t}=X^i(\gamma(t)).
\]
Solving this equation for different initial points produces a family
of maps. The map at parameter value \(t\) sends each initial point to
the point reached after following the field for that parameter interval.

\begin{definition}{Integral curve and local flow}{flow}
An \emph{integral curve} of \(X\) is a smooth curve \(\gamma:I\to M\)
satisfying \(\dot\gamma(t)=X_{\gamma(t)}\).

The \emph{local flow} of \(X\) is a smooth map
\(\Phi:\mathcal D\to M\), where \(\mathcal D\) is an open neighborhood
of \(\{0\}\times M\) in \(\R\times M\), such that
\[
 \Phi_0(p)=p,\qquad
 \frac{\dd}{\dd t}\Phi_t(p)=X_{\Phi_t(p)},
 \qquad \Phi_t(p)=\Phi(t,p).
\]
For each \(p\), the allowed parameter values form an open interval
containing zero. The curve \(t\mapsto\Phi_t(p)\) is the integral
curve starting at \(p\).
\end{definition}

Existence, uniqueness, and smooth dependence on initial conditions
follow from the corresponding ordinary differential-equation results.
Uniqueness gives
\[
 \Phi_{t+s}(p)=\Phi_t(\Phi_s(p)),\qquad
 \Phi_{-t}=\Phi_t^{-1}
\]
on the relevant common domains. Following the field for \(s\) and
then for \(t\) is the same as following it for \(t+s\). Each fixed-time
map is a diffeomorphism between its domain and image.

A local flow need not exist for every real parameter value. For example,
\(X=x^2\partial_x\) on \(\R\) gives
\(\Phi_t(a)=a/(1-at)\), whose denominator vanishes at \(t=1/a\) when
\(a\neq0\). A field is \emph{complete} when every maximal integral
curve is defined for all real times. The derivatives below use only
parameter values near zero.

For the sphere field \(K(x,y,z)=(-y,x,0)\), the flow is the rotation
\[
 \Phi_s(x,y,z)=(x\cos s-y\sin s,\ x\sin s+y\cos s,\ z).
\]
In angular coordinates, \(K=\partial_\varphi\) and the same map is
\((\vartheta,\varphi)\mapsto(\vartheta,\varphi+s)\) wherever that chart
applies. The rotation itself is defined globally, including when a
trajectory crosses the longitude cut. Its fixed points are the poles.

\subsection{Lie differentiation}

Fix \(p\) and let \(q=\Phi_t(p)\). To compare \(Y_q\) with \(Y_p\),
use the inverse flow to bring \(Y_q\) back to \(T_pM\),
\[
 Y_q\in T_qM
 \quad\longmapsto\quad
 \dd(\Phi_{-t})_q(Y_q)\in T_pM.
\]
This is a comparison supplied by the flow of \(X\), not by a connection.
The Lie derivative measures the resulting first-order change.

\begin{definition}{Lie derivative}{lie-derivative}
Let \(\Phi_t\) be the local flow of \(X\). For a vector field \(Y\),
\[
 (\Lie_XY)_p
 =\lim_{t\to0}
 \frac{\dd(\Phi_{-t})_{\Phi_t(p)}(Y_{\Phi_t(p)})-Y_p}{t}.
\]
For a covariant tensor field \(A\),
\[
 \Lie_XA=\left.\frac{\dd}{\dd t}\right|_{t=0}\Phi_t^*A.
\]
For a scalar field, \(\Lie_Xf=X(f)\). On mixed tensors, Lie
differentiation is determined by the tensor product rule and
compatibility with contraction.
\end{definition}

In the vector formula, both terms in the numerator lie in \(T_pM\).
For covariant tensors, pullback gives the corresponding comparison at
the same point. This explains the opposite directions of the two maps:
vectors are brought back with the inverse differential, while covariant
tensors are brought back by pullback.


\begin{theorem}{Lie derivative of a vector field}{lie-bracket-derivative}
For smooth vector fields \(X,Y\),
\[
 \Lie_XY=[X,Y].
\]
The condition \([X,Y]=0\) is equivalent to preservation of \(Y\) by
the local flow of \(X\), meaning
\[
 \dd(\Phi_t)_p(Y_p)=Y_{\Phi_t(p)}.
\]
It is also equivalent to local commutation of the flows of \(X\) and
\(Y\).
\end{theorem}

\begin{proof}
In one chart, the flow and the differential of its inverse have the
first-order expansions
\[
 \begin{aligned}
 \Phi_t^i(x)&=x^i+tX^i(x)+O(t^2),\\
 (\dd\Phi_{-t})^i{}_j\big|_{\Phi_t(x)}
 &=\delta^i{}_j-t\partial_jX^i(x)+O(t^2).
 \end{aligned}
\]
Also,
\(Y^i(\Phi_t(x))=Y^i(x)+tX^j\partial_jY^i(x)+O(t^2)\).
Multiplying gives
\[
 \dd(\Phi_{-t})_{\Phi_t(p)}(Y_{\Phi_t(p)})
 =Y_p+t\bigl(X^j\partial_jY^i-Y^j\partial_jX^i\bigr)_p
      \partial_i|_p+O(t^2).
\]
The defining limit is therefore \eqref{eq:bracket-components}.
The term \(X^j\partial_jY^i\) comes from the change in \(Y\) along
the flow. The term \(-Y^j\partial_jX^i\) comes from bringing its value
back by the inverse differential.

Write \(\Phi_t^*Y\) for the vector field obtained by the same
inverse-differential comparison. The flow law gives
\[
 \frac{\dd}{\dd t}\Phi_t^*Y=\Phi_t^*(\Lie_XY).
\]
Thus \(\Lie_XY=0\) is equivalent to \(\Phi_t^*Y=Y\), which is the
stated preservation condition. When \(\Phi_t\) preserves \(Y\), it
maps each integral curve of \(Y\) to an integral curve of \(Y\).
Uniqueness of integral curves makes the two flows commute locally.
Conversely, differentiating their commutation in the flow parameter of
\(Y\) gives the preservation condition.
\end{proof}

\begin{examplebox}{The order of two flows}{flow-commutator}
On \(\R^2\), let \(X=\partial_x\) and \(Y=x\partial_y\). Their flows
are \(\Phi_t(x,y)=(x+t,y)\) and \(\Psi_s(x,y)=(x,y+sx)\). The two
orders give
\[
 \Psi_s\circ\Phi_t(x,y)=(x+t,y+sx+st),\qquad
 \Phi_t\circ\Psi_s(x,y)=(x+t,y+sx).
\]
They differ in the \(y\)-direction by \(st\). Equivalently,
\[
 \Psi_{-s}\circ\Phi_{-t}\circ\Psi_s\circ\Phi_t(x,y)=(x,y+st).
\]
The residual displacement agrees with \([X,Y]=\partial_y\). The
bracket detects the failure of the two flows to commute.
\end{examplebox}

When frame fields have commuting flows, their flow parameters can serve
as independent coordinates. This gives a criterion for recognizing
coordinate frames.

\begin{proposition}{When a frame comes from coordinates}{coordinate-frame}
Smooth frame fields \(e_1,\ldots,e_n\) form a coordinate frame locally
around each point if and only if \([e_i,e_j]=0\) for all \(i,j\).
\end{proposition}

\begin{proof}
Coordinate partial derivatives commute. Conversely, suppose the frame
fields commute, with local flows \(\Phi_i^t\). Fix a point \(p\) and
set
\[
 F(t^1,\ldots,t^n)=\Phi_1^{t^1}\circ\cdots\circ\Phi_n^{t^n}(p)
\]
for sufficiently small parameters. The differential at zero sends the
standard basis to \(e_1|_p,\ldots,e_n|_p\), so it is invertible.
The inverse function theorem makes \(F\) a local parametrization.
Commutation allows each flow factor to be differentiated independently,
giving \(\dd F(\partial/\partial t^i)=e_i\). Hence the inverse map
defines coordinates with the required coordinate frame.
\end{proof}

Preservation of \(Y\) does not mean that \(Y(f)\) is constant along
\(X\) for every function \(f\). For instance, \(\partial_x\) and
\(\partial_y\) commute, but \(\partial_y(xy)=x\) varies along
\(\partial_x\). The bracket condition says that the two orders of
differentiation agree, not that their common value is zero.

\subsection{Lie derivatives of differential forms}

For forms, the Lie derivative can be computed without solving for the
flow explicitly. The required additional operation is insertion of the
vector field into a form.

For a \(k\)-form \(\alpha\), \(k\geq1\), the \emph{interior product}
\(\iota_X\alpha\) inserts \(X\) into the first argument,
\[
 (\iota_X\alpha)(X_2,\ldots,X_k)=\alpha(X,X_2,\ldots,X_k).
\]
It is a \((k-1)\)-form. For a one-form, \(\iota_X\alpha=\alpha(X)\)
is a function. Set \(\iota_Xf=0\) for functions. The wedge-product rule
is
\[
 \iota_X(\alpha\wedge\beta)
 =\iota_X\alpha\wedge\beta+(-1)^k\alpha\wedge\iota_X\beta,
 \qquad \alpha\in\Omega^k(M).
\]

\begin{theorem}{Cartan's formula}{cartan}
For every differential form \(\alpha\),
\begin{equation}
 \Lie_X\alpha=\iota_X\dd\alpha+\dd(\iota_X\alpha).
 \label{eq:cartan}
\end{equation}
\end{theorem}

\begin{proof}
Both sides preserve degree. On the right, one operation raises the
degree and the other lowers it. Expanding their graded product rules
shows that the mixed terms cancel, leaving the ordinary product rule
for wedge products. On functions, the right-hand side is
\(\iota_X\dd f=X(f)\).
Naturality of \(\dd\) under pullback gives
\(\Lie_X\dd f=\dd(X(f))\), which also agrees with the right-hand side.
Locally, every form is a sum of products of functions and coordinate
differentials. Agreement on these factors and the product rule therefore
imply agreement on every form.
\end{proof}

Lie differentiation of the scalar pairing gives
\[
 X\bigl(\alpha(Y)\bigr)=(\Lie_X\alpha)(Y)+\alpha([X,Y]).
\]
For \(\alpha=\alpha_i\dd x^i\), the product rule and
\(\Lie_X\dd x^i=\dd(X^i)\) give
\begin{equation}
 (\Lie_X\alpha)_i=X^j\partial_j\alpha_i+\alpha_j\partial_iX^j.
 \label{eq:lie-one-form}
\end{equation}
The second term is needed because the flow also acts on the covector
basis. Lie differentiation preserves form degree, while exterior
differentiation raises it. The two operations commute,
\(\Lie_X\dd\alpha=\dd\Lie_X\alpha\).

\subsection{Exterior differentiation in a general frame}

The coordinate formula for \(\dd\alpha\) uses commuting coordinate
fields. For a one-form \(\alpha\) and arbitrary vector fields \(X,Y\),
the corresponding formula includes a bracket term,
\begin{equation}
 (\dd\alpha)(X,Y)
 =X\bigl(\alpha(Y)\bigr)-Y\bigl(\alpha(X)\bigr)-\alpha([X,Y]).
 \label{eq:intrinsic-d}
\end{equation}
Expanding in coordinates shows why: derivatives of the components of
\(X\) and \(Y\) cancel against the bracket term, leaving
\((\partial_i\alpha_j-\partial_j\alpha_i)X^iY^j\).

Let \(\varepsilon^1,\ldots,\varepsilon^n\) be dual to a local frame
\(e_1,\ldots,e_n\), and define its commutator coefficients by
\([e_b,e_c]=C^a{}_{bc}e_a\). Since
\(\varepsilon^a(e_b)=\delta^a{}_b\) is constant, the first two
terms of \eqref{eq:intrinsic-d} vanish, leaving
\((\dd\varepsilon^a)(e_b,e_c)=-C^a{}_{bc}\). Expanding in the coframe
gives
\begin{equation}
 \dd\varepsilon^a=-\frac12 C^a{}_{bc}\,
                         \varepsilon^b\wedge\varepsilon^c.
 \label{eq:coframe-bracket}
\end{equation}
The functions \(C^a{}_{bc}\) depend on the chosen frame and do not
transform as tensor components under an arbitrary varying frame change.
For a coordinate coframe they vanish, giving \(\dd(\dd x^i)=0\).
A general coframe need not consist of differentials of coordinates.

\begin{remarkbox}[Lie differentiation and covariant differentiation]
Lie differentiation compares a field using the flow of \(X\).
Covariant differentiation compares it using a chosen connection.
These rules need not agree. For the flat connection on \(\R^2\), take
\(X=x\partial_x+y\partial_y\) and \(Y=\partial_x\). Then
\[
 \nabla_XY=0,\qquad \Lie_XY=[X,Y]=-Y.
\]
The flow of \(X\) is a dilation and rescales tangent vectors. Its
comparison differs from parallel transport for the flat connection,
which preserves the constant Cartesian field. Moreover,
\(\Lie_XY\) is not linear over \(\Cinf(M)\) in \(X\), whereas
\(\nabla_XY\) must be.
\end{remarkbox}

\begin{studentwork}[3.8cm]{rotation and the sphere one-form}
For the sphere rotation field \(K=\partial_\varphi\) and the one-form
\(\beta=\iota^*(x\,\dd y-y\,\dd x)\), compute
\(\iota_K\beta\) and \(\iota_K\dd\beta\). Show that
\(\Lie_K\beta=0\) using both Cartan's formula and the rotation flow.
\end{studentwork}

\section{Metrics and measurement}
\label{sec:metrics}
\slideref{18}

A smooth manifold supplies tangent vectors and covectors, but it does not
assign scalar products to tangent vectors. A metric supplies a bilinear
form in each tangent space, with smooth dependence on the base point.
This is the pointwise generalization of the Minkowski scalar product
used in SR1.

\subsection{Metrics and signature}

\begin{definition}{Pseudo-Riemannian metric}{metric}
A \emph{pseudo-Riemannian metric} is a smooth tensor field \(g\) of
type \((0,2)\) whose value at each \(p\in M\) is a symmetric,
nondegenerate bilinear form
\[
 g_p:T_pM\times T_pM\longrightarrow\R.
\]
Nondegeneracy means
\[
 \bigl[g_p(v,w)=0\text{ for every }w\in T_pM\bigr]
 \quad\Longrightarrow\quad v=0.
\]
The pair \((M,g)\) is a \emph{pseudo-Riemannian manifold}.
\end{definition}

Both input vectors belong to the same tangent space. In coordinates,
\begin{equation}
 g=g_{ij}\,\dd x^i\otimes\dd x^j,
 \qquad g_{ij}=g(\partial_i,\partial_j).
 \label{eq:metric-tensor}
\end{equation}
The component matrix is symmetric and invertible. Smoothness means that
each coefficient \(g_{ij}\) is smooth. On a chart overlap,
\[
 g'_{ab}=\frac{\partial x^i}{\partial y^a}
         \frac{\partial x^j}{\partial y^b}g_{ij}.
\]
The components change, but their evaluation on the same two vectors
gives the same scalar.

Every nondegenerate real symmetric bilinear form has a basis in which
its matrix is \(\operatorname{diag}(I_{n_+},-I_{n_-})\), with
\(n_++n_-=n\). The numbers of positive and negative entries are
independent of the basis. Their pair \((n_+,n_-)\) is the
\emph{signature}. For a smooth nondegenerate metric, the signature is
constant on each connected component. A change of sign of an eigenvalue
would require it to pass through zero, contradicting nondegeneracy.

A positive-definite metric has signature \((n,0)\) and is called
\emph{Riemannian}. A \emph{Lorentzian} metric has one sign different
from the others. In four dimensions we use \((n_+,n_-)=(1,3)\), or
\((+,-,-,-)\).

For a Riemannian metric, lengths and angles at a point are
\[
 \|v\|_g=\sqrt{g(v,v)},\qquad
 \cos\theta=\frac{g(v,w)}{\|v\|_g\|w\|_g}
 \quad (v,w\neq0).
\]
The length of a piecewise smooth curve is
\[
 L_g(\gamma)=\int_a^b
 \sqrt{g_{\gamma(t)}(\dot\gamma(t),\dot\gamma(t))}\,\dd t.
\]
Under a regular reparametrization, the integrand acquires the factor
\(|\dd t/\dd s|\), so the integral is unchanged. On a connected
Riemannian manifold, the distance between two points is the infimum
of the lengths of curves joining them. It is not generally the length
of an arbitrarily chosen joining curve.


\begin{remarkbox}[Nondegeneracy and null vectors]
A Lorentzian metric can have \(g(v,v)=0\) for \(v\neq0\).
Nondegeneracy excludes a nonzero vector orthogonal to \emph{every}
vector, not one orthogonal to itself. Thus \(\sqrt{|g(v,v)|}\) is
not a norm, and the Riemannian angle formula does not apply to
arbitrary Lorentzian vectors.

With our signature, a nonzero vector is timelike, null, or spacelike
according as \(g(v,v)>0\), \(g(v,v)=0\), or \(g(v,v)<0\).
Distinguishing future from past also requires a consistent choice
of time orientation.
\end{remarkbox}

\subsection{Metric duality and the inverse metric}

Fix \(v\in T_pM\) and allow the second argument of \(g_p(v,w)\) to
vary. The result is a covector, denoted by \(v^\flat\),
\[
 v^\flat(w)=g_p(v,w),\qquad
 \flat_p:T_pM\longrightarrow T_p^*M.
\]
If \(v^\flat=0\), nondegeneracy gives \(v=0\). Since the two spaces
have the same finite dimension, \(\flat_p\) is an isomorphism. Its
inverse is denoted by \(\sharp_p:T_p^*M\to T_pM\). These are the
\emph{musical isomorphisms}.

Let \((g^{ij})\) be the inverse of \((g_{ij})\), so that
\(g^{ik}g_{kj}=\delta^i{}_j\). Then
\begin{equation}
 v^\flat=g_{ij}v^j\dd x^i,
 \qquad \alpha^\sharp=g^{ij}\alpha_j\partial_i.
 \label{eq:metric-duality}
\end{equation}
These are index lowering and raising. The maps vary smoothly with the
point and therefore also act on fields. Unlike evaluation \(\alpha(v)\),
they depend on the metric.

The inverse metric is a tensor of type \((2,0)\),
\[
 g^{-1}=g^{ij}\partial_i\otimes\partial_j,
 \qquad
 g^{-1}(\alpha,\beta)=g(\alpha^\sharp,\beta^\sharp)
                     =g^{ij}\alpha_i\beta_j.
\]
It acts on covectors rather than vectors. Even when the matrices
\((g_{ij})\) and \((g^{ij})\) have the same numerical entries, they
represent tensors of different types.

The differential \(\dd f\) gives the rate of change of \(f\) in any
direction. The metric lets us express this same linear measurement as a
scalar product with a vector. That vector field is the \emph{metric
gradient}, obtained by raising the index of \(\dd f\),
\begin{equation}
 \grad_g f=(\dd f)^\sharp
          =g^{ij}(\partial_j f)\partial_i,
 \qquad g(\grad_g f,X)=X(f).
 \label{eq:gradient}
\end{equation}
The one-form \(\dd f\) is defined by the smooth structure alone.
The vector \(\grad_g f\) is defined only after choosing \(g\).

\section{Sphere and spacetime metrics}
\label{sec:examples}
\slideref{19}

\subsection{The round metric on the sphere}
\label{sec:sphere-metric}

The embedding \(\iota:S^2\hookrightarrow\R^3\) lets us measure tangent
vectors using the Euclidean metric \(\delta\) of the ambient space.
Its pullback is the \emph{induced metric},
\[
 g_{S^2}=\iota^*\delta,\qquad
 (g_{S^2})_p(v,w)=\delta_{\iota(p)}(\dd\iota_pv,\dd\iota_pw).
\]
The differential \(\dd\iota_p\) is injective, so a nonzero tangent
vector has a nonzero ambient representative. Its Euclidean square is
positive, making the induced metric Riemannian.

Using the coordinate vectors from Section~\ref{sec:sphere-tangent},
\[
 g_{\vartheta\vartheta}=e_\vartheta\cdot e_\vartheta=1,
 \qquad g_{\vartheta\varphi}=0,
 \qquad g_{\varphi\varphi}=e_\varphi\cdot e_\varphi=\sin^2\vartheta.
\]
Thus
\begin{equation}
 g_{S^2}=\dd\vartheta\otimes\dd\vartheta
          +\sin^2\vartheta\,\dd\varphi\otimes\dd\varphi.
 \label{eq:sphere-metric}
\end{equation}
The factor \(\sin^2\vartheta\) reflects the smaller radius of a
latitude circle, since \(\|e_\varphi\|=\sin\vartheta\).
The component order is \((\vartheta,\varphi)\). Reversing that order
gives \(\operatorname{diag}(\sin^2\vartheta,1)\), without changing
the metric. For a sphere of radius \(R\), the angular metric is
multiplied by \(R^2\).

On this angular chart, the inverse metric gives
\[
 \grad f=(\partial_\vartheta f)\partial_\vartheta
          +\csc^2\vartheta(\partial_\varphi f)\partial_\varphi.
\]
The formula does not apply at the poles because the angular chart does
not include them. The apparent singularity of its coefficients does
not imply that the metric or a smooth gradient is singular there.

Rescaling the coordinate vectors gives the orthonormal frame and its
dual coframe,
\[
 E_1=\partial_\vartheta,\qquad
 E_2=\csc\vartheta\,\partial_\varphi,
 \qquad
 \varepsilon^1=\dd\vartheta,\qquad
 \varepsilon^2=\sin\vartheta\,\dd\varphi.
\]
In this frame,
\[
 g_{S^2}=\varepsilon^1\otimes\varepsilon^1+\varepsilon^2\otimes\varepsilon^2.
\]
The metric components are constant, but the frame is not a coordinate
frame. Indeed,
\[
 [E_1,E_2]=-\cot\vartheta\,E_2,\qquad
 \dd\varepsilon^1=0,\qquad
 \dd\varepsilon^2=\cot\vartheta\,\varepsilon^1\wedge\varepsilon^2.
\]
The bracket vanishes on the equator, but not on any nonempty open
subset of the angular chart. By the coordinate-frame criterion, no
coordinates on such a set have \(E_1,E_2\) as their coordinate vectors.
Constant metric components in a varying noncoordinate frame therefore
do not establish flatness.

The rotation field and the one-form introduced earlier are now related,
\[
 K=\partial_\varphi,\qquad K^\flat=\sin^2\vartheta\,\dd\varphi=\beta,
 \qquad g(K,K)=\sin^2\vartheta.
\]
Their definitions did not require a metric. The identification
\(K^\flat=\beta\) uses the particular round metric.

The stereographic chart is regular at the south pole. Differentiating
its inverse parametrization gives
\begin{equation}
 g_{S^2}=\frac{4}{(1+u^2+v^2)^2}
           (\dd u\otimes\dd u+\dd v\otimes\dd v).
 \label{eq:stereographic-metric}
\end{equation}
At \((u,v)=(0,0)\), the matrix is \(4I_2\). This describes the same
metric in coordinates that remain valid where the angular chart fails.


\begin{studentwork}[3.2cm]{the sphere metric in two frames}
Verify the orthonormal frame and coframe identities above, then
derive \eqref{eq:stereographic-metric}.
\end{studentwork}

\begin{studentwork}[3.3cm]{the gradient of height}
For \(f=z\) on \(S^2\), compute \(\dd f\), \(\grad f\), and
\(g(\grad f,\grad f)\). Find the gradient at the poles using
its ambient tangent-vector representation.
\end{studentwork}

\subsection{Lorentzian line elements and proper time}

In inertial coordinates \(x^\mu=(ct,x,y,z)\), the Minkowski metric is
\[
 \eta=\eta_{\mu\nu}\,\dd x^\mu\otimes\dd x^\nu,
 \qquad (\eta_{\mu\nu})=\operatorname{diag}(1,-1,-1,-1).
\]
A general spacetime metric is written
\(g=g_{\mu\nu}(x)\,\dd x^\mu\otimes\dd x^\nu\). Its
\emph{line element} is the quadratic-form notation
\begin{equation}
 \dd s^2=g_{\mu\nu}(x)\,\dd x^\mu\dd x^\nu.
 \label{eq:line-element}
\end{equation}
Along a curve \(\gamma(\lambda)\), this means
\[
 \dd s^2=g_{\gamma(\lambda)}(\dot\gamma,\dot\gamma)\,\dd\lambda^2,
 \qquad \dot\gamma=\frac{\dd\gamma}{\dd\lambda}.
\]
The products \(\dd x^\mu\dd x^\nu\) in a line element are not wedge
products. A symmetric tensor cannot be represented by
\(g_{\mu\nu}\dd x^\mu\wedge\dd x^\nu\), since that expression
vanishes by antisymmetry.

In affine Minkowski coordinates, the finite displacement between two
events has invariant square
\(\Delta s^2=\eta_{\mu\nu}\Delta x^\mu\Delta x^\nu\).
On a general manifold, coordinate differences do not form the
components of an intrinsic displacement vector. Therefore,
\(g_{\mu\nu}(p)\Delta x^\mu\Delta x^\nu\) is not an exact invariant
separation of arbitrary events. The line element instead evaluates
the metric on a tangent vector at one point of a curve.

For a future-directed timelike worldline in a time-oriented spacetime,
the elapsed proper time is
\begin{equation}
 \tau[\gamma]=\frac1c\int_a^b
 \sqrt{g_{\gamma(\lambda)}(\dot\gamma(\lambda),\dot\gamma(\lambda))}
 \,\dd\lambda.
 \label{eq:proper-time}
\end{equation}
Here \(\dd s^2\) has units of length squared. The integral depends on
the worldline, not only on its endpoints. With proper time as the
parameter, the four-velocity \(U=\dd\gamma/\dd\tau\) satisfies
\(g(U,U)=c^2\), or \(g(U,U)=1\) in units \(c=1\).

\subsection{Local Lorentz frames}
\label{sec:local-lorentz}

At any event \(p\), the Lorentzian metric admits a basis
\(e_0|_p,\ldots,e_3|_p\) with \(g_p(e_a,e_b)=\eta_{ab}\).
An orthonormal frame can also be chosen smoothly on a sufficiently
small neighborhood. Extend the basis at \(p\) to local fields and
orthonormalize using the metric. After shrinking the neighborhood,
the normalization factors remain nonzero, so the construction is smooth.

Write the frame and its dual coframe as
\[
 e_a=e_a{}^\mu\partial_\mu,\qquad
 \varepsilon^a=\varepsilon^a{}_\mu\dd x^\mu,
 \qquad \varepsilon^a(e_b)=\delta^a{}_b.
\]
Then
\[
 g=\eta_{ab}\,\varepsilon^a\otimes\varepsilon^b,
 \qquad g_{\mu\nu}=\eta_{ab}\varepsilon^a{}_\mu\varepsilon^b{}_\nu.
\]
Here \(a,b=0,1,2,3\) label the orthonormal frame, while \(\mu,\nu\)
label the coordinate basis. The frame need not be a coordinate frame,
just as in the sphere example.

To retain the component convention of SR1, write another orthonormal
frame as
\[
 e'_a=e_b(\Lambda^{-1})^b{}_a,
 \qquad \Lambda^{\mathsf T}\eta\Lambda=\eta.
\]
The components of the same vector then satisfy
\(v'^a=\Lambda^a{}_b v^b\). Thus Lorentz transformations relate
orthonormal bases of a tangent space. A change of spacetime coordinates
has a general invertible Jacobian and need not relate orthonormal
coordinate bases. These are different uses of a basis transformation.

A linear change of a chart can make its coordinate basis agree with
a prescribed basis at \(p\). In particular, one can arrange
\(g_{\mu\nu}(p)=\eta_{\mu\nu}\). This fixes the metric components
at one event. It does not make them constant on a neighborhood or
make their first derivatives vanish. Appendix~\ref{sec:inertial-point}
constructs the stronger local inertial coordinate choice.

\subsection{Metric invariance under a flow}

A diffeomorphism \(F:M\to M\) is an \emph{isometry} when \(F^*g=g\).
Equivalently,
\[
 g_{F(p)}(\dd F_pv,\dd F_pw)=g_p(v,w)
\]
for all \(p,v,w\). The map sends points and vectors to new ones while
preserving their metric pairings. This is an additional property of
the map, not an automatic consequence of changing coordinates.

The Lie derivative tests whether a flow consists of local isometries.
For the metric, the tensor product rule gives
\begin{equation}
 (\Lie_Xg)_{ij}
 =X^k\partial_k g_{ij}+g_{kj}\partial_iX^k+g_{ik}\partial_jX^k.
 \label{eq:lie-metric}
\end{equation}
The first term measures how the metric coefficients vary along \(X\).
The other two account for the flow's action on the two vector arguments.
The condition \(\Lie_Xg=0\) is equivalent to \(\Phi_t^*g=g\) on the
flow domains. Such an \(X\) is called a \emph{Killing vector field}.
This definition requires a metric and a vector field, but no connection.

For the round sphere, \(K=\partial_\varphi\) has constant components
in the angular chart and the metric coefficients are independent of
\(\varphi\). Hence \(\Lie_Kg=0\). Its flow is the rotation already
constructed, which preserves the induced metric globally. Since
\(\beta=g(K,\cdot)\), the product rule also gives
\[
 \Lie_K\beta=(\Lie_Kg)(K,\cdot)+g(\Lie_KK,\cdot)=0.
\]
This relates the earlier invariance of the sphere one-form to the
symmetry of the metric.

\subsection{From the metric to covariant geometry}
\slideref{20}

A smooth manifold has local charts and a tangent space at each point.
Charts give components, while transition maps relate the descriptions
of the same tensors and forms. Exterior differentiation uses only the
smooth structure. Lie differentiation compares fields through a local
flow, while a connection defines covariant differentiation and
parallel transport along curves. A metric supplies a nondegenerate
bilinear pairing within each tangent space.

To transport vectors along a curve while preserving their metric
pairing, choose a metric-compatible connection. Requiring zero torsion
as well selects the Levi--Civita connection.
Appendix~\ref{sec:levi-civita} derives this connection and completes
the sphere transport calculation. Appendix~\ref{sec:boost-app}
applies these constructions to a noninertial chart on Minkowski
spacetime.

\appendix
\renewcommand{\thesection}{\Alph{section}}

\section{The Levi--Civita connection}
\label{sec:levi-civita}

An affine connection and a metric are independent choices in general.
For a given metric, requiring compatibility with its bilinear forms and
requiring vanishing torsion selects a unique connection. To state these
conditions precisely, first extend covariant differentiation from vectors
to covectors and tensor fields.

\subsection{Covectors and tensor fields}

For a one-form \(\alpha\) and a vector field \(Y\), the pairing
\(\alpha(Y)\) is a scalar function. We require its derivative to obey
\[
 X(\alpha(Y))=(\nabla_X\alpha)(Y)+\alpha(\nabla_XY).
\]
Solving for the first term defines covariant differentiation of covectors,
\[
 (\nabla_X\alpha)(Y)
 =X\bigl(\alpha(Y)\bigr)-\alpha(\nabla_XY).
\]
This is \(\Cinf(M)\)-linear in \(Y\), so it is a one-form. In coordinates,
\begin{equation}
 \nabla_i\alpha_j:=
 (\nabla_{\partial_i}\alpha)(\partial_j)
 =\partial_i\alpha_j-\Gamma^k{}_{ij}\alpha_k.
 \label{eq:covariant-covector}
\end{equation}
The opposite signs of the connection terms ensure the ordinary product rule
for the scalar contraction \(\alpha_jY^j\).

For general tensors, require the Leibniz rule for tensor products and
compatibility with contractions. There is one positive connection term for
each upper index and one negative term for each lower index. For example,
\[
 (\nabla_i T)^a{}_b
 =\partial_iT^a{}_b+\Gamma^a{}_{ic}T^c{}_b
  -\Gamma^c{}_{ib}T^a{}_c.
\]
For a scalar field, \(\nabla_i f=\partial_i f\). These first derivatives
are the components of the one-form \(\dd f\). Differentiating once more
therefore uses the covector rule, giving the \emph{covariant Hessian},
\begin{equation}
 (\nabla\dd f)_{ij}
 =\partial_i\partial_j f-\Gamma^k{}_{ij}\partial_k f.
 \label{eq:hessian}
\end{equation}
Unlike the uncorrected second partial derivatives, this is a covariant
two-tensor. This tensor is also denoted by \(\operatorname{Hess}_\nabla f\).
Its symmetry will depend on the torsion of the connection.


\begin{studentwork}[3.6cm]{the difference of two connections}
Prove that the difference of two affine connections defines a tensor
of type \((1,2)\). Verify the result from the transformation law of
their connection coefficients.
\end{studentwork}

\subsection{Torsion and metric compatibility}

The first condition relates the antisymmetric part of covariant
differentiation to the Lie bracket. The second controls how scalar
products change under differentiation.

\begin{definition}{Torsion}{torsion}
The \emph{torsion} of an affine connection \(\nabla\) is
\[
 T(X,Y)=\nabla_XY-\nabla_YX-[X,Y].
\]
The connection is \emph{torsion free} if \(T(X,Y)=0\) for all vector
fields \(X,Y\).
\end{definition}

The derivative terms from the Leibniz rules cancel, so \(T\) is
\(\Cinf(M)\)-bilinear and defines a tensor of type \((1,2)\). In a
coordinate frame, \([\partial_i,\partial_j]=0\), and therefore
\[
 T^k{}_{ij}=\Gamma^k{}_{ij}-\Gamma^k{}_{ji}.
\]
Consequently, torsion vanishes exactly when the lower indices of the
connection coefficients are symmetric in every coordinate frame. For a general frame \(e_a\), define
\(\nabla_{e_a}e_b=\Gamma^c{}_{ab}e_c\) and
\([e_a,e_b]=C^c{}_{ab}e_c\). Then
\[
 T^c{}_{ab}=\Gamma^c{}_{ab}-\Gamma^c{}_{ba}-C^c{}_{ab}.
\]
The bracket term must be included when the frame is not a coordinate frame.

\begin{definition}{Metric-compatible connection}{metric-compatible}
An affine connection is \emph{compatible with \(g\)} if
\[
 X\bigl(g(Y,Z)\bigr)
 =g(\nabla_XY,Z)+g(Y,\nabla_XZ)
\]
for all vector fields \(X,Y,Z\). Equivalently, the covariant derivative
of the metric vanishes, \(\nabla g=0\).
\end{definition}

In coordinates, compatibility reads
\begin{equation}
 \label{eq:metric-compatible-components}
 \partial_i g_{jk}
 =\Gamma^\ell{}_{ij}g_{\ell k}
  +\Gamma^\ell{}_{ik}g_{j\ell}.
\end{equation}
This condition links differentiation of the metric components to the
connection, rather than requiring the components to be constant.

\subsection{Existence and uniqueness}

\begin{theorem}{Levi--Civita connection}{levi-civita}
Every pseudo-Riemannian metric \(g\) determines a unique torsion-free,
metric-compatible affine connection. This is its \emph{Levi--Civita
connection}.
\end{theorem}

To determine \(\nabla_XY\), it is enough to determine its scalar
product with every vector field \(Z\). Nondegeneracy then fixes the
vector itself. Combining metric compatibility with zero torsion gives
the \emph{Koszul formula},
\begin{equation}
 \label{eq:koszul}
 \begin{aligned}
 2g(\nabla_XY,Z)
 ={}&X\bigl(g(Y,Z)\bigr)+Y\bigl(g(Z,X)\bigr)-Z\bigl(g(X,Y)\bigr)\\
    &-g(X,[Y,Z])+g(Y,[Z,X])+g(Z,[X,Y]).
 \end{aligned}
\end{equation}
To derive it, expand the three differentiated scalar products using
compatibility. Their signs are chosen so that the unwanted derivatives
cancel after using \(\nabla_XY-\nabla_YX=[X,Y]\), together with its
cyclic versions. The remaining derivative is
\(2g(\nabla_XY,Z)\), and the interchanges produce the three bracket
terms displayed above.

The right-hand side depends only on \(g,X,Y,Z\), so nondegeneracy gives
uniqueness. For existence, use the same formula as a definition and
check the required properties. Denote half of its right-hand side by
\(B(X,Y,Z)\). The bracket and product rules give
\(B(X,Y,fZ)=fB(X,Y,Z)\). In this calculation the terms containing
\(X(f)g(Y,Z)\) and \(Y(f)g(X,Z)\) cancel in pairs. Hence \(B\) defines a
smooth one-form in \(Z\) for fixed \(X,Y\). Applying \(\sharp\) to
this one-form defines a unique smooth field \(\nabla_XY\).

It remains to check that this field defines a connection and has the two
promised properties. The same rules give
\[
 B(fX,Y,Z)=fB(X,Y,Z),\qquad
 B(X,fY,Z)=X(f)g(Y,Z)+fB(X,Y,Z).
\]
Together with real bilinearity, these are the connection axioms.
The sum and difference identities are
\[
 \begin{aligned}
 B(X,Y,Z)+B(X,Z,Y)&=X\bigl(g(Y,Z)\bigr),\\
 B(X,Y,Z)-B(Y,X,Z)&=g([X,Y],Z).
 \end{aligned}
\]
The first is metric compatibility. By nondegeneracy, the second gives
\(\nabla_XY-\nabla_YX=[X,Y]\). This proves existence as well as uniqueness,
without first choosing coordinates.

\begin{proposition}{Levi--Civita coefficients}{christoffel-metric}
In a coordinate frame, the Levi--Civita connection has coefficients
\begin{equation}
 \label{eq:christoffel-metric}
 \Gamma^k{}_{ij}
 =\frac12g^{k\ell}
 \left(\partial_i g_{j\ell}+\partial_j g_{i\ell}
       -\partial_\ell g_{ij}\right).
\end{equation}
\end{proposition}

To see the coordinate calculation directly, write the three metric
compatibility equations and use the symmetry of the lower connection indices.
Adding the equations with differentiation directions \(i,j\) and subtracting
the equation with direction \(\ell\) gives
\[
 \partial_i g_{j\ell}+\partial_j g_{i\ell}-\partial_\ell g_{ij}
 =2g_{k\ell}\Gamma^k{}_{ij}.
\]
Multiplication by \(g^{m\ell}\) gives the stated formula. Equivalently,
evaluate the Koszul formula on coordinate fields, whose brackets vanish.
Formula~\eqref{eq:christoffel-metric} is specific to the Levi--Civita
connection. A general affine connection is not determined by a metric.

\Needspace{6.5cm}
\begin{studentwork}[4.4cm]{constructing the metric connection}
Derive the Koszul formula from metric compatibility and zero torsion.
Verify the linearity and product-rule identities for \(B\) used in the
existence proof.
\end{studentwork}

\subsection{Parallel transport and scalar products}

For a metric-compatible connection and fields \(V,W\) along a curve,
\begin{equation}
 \label{eq:metric-transport}
 \frac{\dd}{\dd t}g_{\gamma(t)}(V,W)
 =g_{\gamma(t)}\left(\frac{DV}{\dd t},W\right)
  +g_{\gamma(t)}\left(V,\frac{DW}{\dd t}\right).
\end{equation}
The coordinate formula follows directly from
\eqref{eq:metric-compatible-components}, and is valid even when the fields
are defined only along the curve. If both fields are parallel, the
right-hand side vanishes. Thus parallel transport is a linear isometry
between the endpoint tangent spaces,
\[
 g_{\gamma(b)}(P_\gamma v,P_\gamma w)
 =g_{\gamma(a)}(v,w).
\]
It preserves Riemannian lengths and angles, and preserves the causal type of
nonzero vectors in Lorentzian geometry. No positive-definite norm is needed
for the latter statement.

Torsion also expresses the relation between the two derivatives of vector
fields,
\[
 \Lie_XY=[X,Y]=\nabla_XY-\nabla_YX-T(X,Y).
\]
For the Levi--Civita connection, the torsion term is absent, but the term
\(\nabla_YX\) remains. Lie invariance and parallelism are therefore still
different conditions.

The Hessian in \eqref{eq:hessian} is symmetric for a torsion-free
connection, since the lower connection indices are symmetric. It need
not vanish. By contrast, \(\dd^2f\) always vanishes. The covariant
Hessian and the second exterior derivative are different constructions.

For a one-form \(\alpha\), combine \eqref{eq:intrinsic-d} with the
covector product rule to obtain
\begin{equation}
 (\dd\alpha)(X,Y)=(\nabla_X\alpha)(Y)-(\nabla_Y\alpha)(X)
                    +\alpha(T(X,Y)).
 \label{eq:d-and-connection}
\end{equation}
Thus, for a torsion-free connection,
\((\dd\alpha)_{ij}=\nabla_i\alpha_j-\nabla_j\alpha_i\).
The symmetric connection coefficients cancel under antisymmetrization.
This equality expresses an already intrinsic exterior derivative in terms
of a connection. It does not make \(\dd\) dependent on that connection.

\Needspace{5.5cm}
\begin{studentwork}[3.0cm]{torsion and the Hessian}
Show that for any affine connection,
\[
 \operatorname{Hess}_\nabla f(X,Y)-\operatorname{Hess}_\nabla f(Y,X)
 =-\dd f\bigl(T(X,Y)\bigr).
\]
\end{studentwork}

\Needspace{5.8cm}
\begin{studentwork}[3.2cm]{metric symmetry}
For the Levi--Civita connection, show that
\[
 (\Lie_Xg)(Y,Z)=g(\nabla_YX,Z)+g(Y,\nabla_ZX).
\]
Conclude that a Killing field satisfies \(\nabla_iX_j+\nabla_jX_i=0\).
\end{studentwork}

\subsection{The sphere connection and transport around a latitude}
\label{sec:sphere-transport}

For the round metric with line element
\(\dd s^2=(\dd\vartheta)^2+\sin^2\vartheta\,(\dd\varphi)^2\), the only nonzero
coordinate coefficients are
\begin{equation}
 \Gamma^\vartheta{}_{\varphi\varphi}=-\sin\vartheta\cos\vartheta,
 \qquad
 \Gamma^\varphi{}_{\vartheta\varphi}
 =\Gamma^\varphi{}_{\varphi\vartheta}=\cot\vartheta.
 \label{eq:sphere-gamma}
\end{equation}
For example,
\(\Gamma^\vartheta{}_{\varphi\varphi}
=-\tfrac12\partial_\vartheta(\sin^2\vartheta)\), while
\(\Gamma^\varphi{}_{\vartheta\varphi}
=\tfrac12\csc^2\vartheta\,\partial_\vartheta(\sin^2\vartheta)\).
All other terms vanish because \(g_{\vartheta\vartheta}=1\),
\(g_{\vartheta\varphi}=0\), and the coefficients are independent of
\(\varphi\).

The same result follows without the Christoffel formula. For the embedded
coordinate vectors, write the normal vector as \(\sigma\). Direct
differentiation gives
\[
 \partial_\vartheta e_\vartheta=-\sigma,\qquad
 \partial_\vartheta e_\varphi=\partial_\varphi e_\vartheta
 =\cot\vartheta\,e_\varphi,
\]
\[
 \partial_\varphi e_\varphi
 =-\sin\vartheta\cos\vartheta\,e_\vartheta
  -\sin^2\vartheta\,\sigma.
\]
Projecting away the normal terms gives \eqref{eq:sphere-gamma}. Tangential
projection is metric compatible because normal and tangent vectors are
orthogonal. It is torsion free because antisymmetrized ambient derivatives
of tangent fields give their intrinsic Lie bracket. Uniqueness therefore
identifies this projection connection with the Levi--Civita connection.

Now follow the latitude \(\vartheta=\vartheta_0\), with
\(0<\vartheta_0<\pi\), using \(\varphi\) as the parameter. Along this
circle, write a tangent field in the orthonormal frame as
\[
 V(\varphi)=a(\varphi)E_1+b(\varphi)E_2,
 \qquad E_1=\partial_\vartheta,\quad
 E_2=\csc\vartheta\,\partial_\varphi.
\]
The connection coefficients give
\[
 \nabla_{\partial_\varphi}E_1=\cos\vartheta_0\,E_2,\qquad
 \nabla_{\partial_\varphi}E_2=-\cos\vartheta_0\,E_1.
\]
Consequently, the parallel-transport equation is
\[
 a'=\cos\vartheta_0\,b,\qquad b'=-\cos\vartheta_0\,a.
\]
These equations rotate the component pair at constant angular rate
\(-\cos\vartheta_0\). Writing \(c_0=\cos\vartheta_0\) and
\(\Delta\varphi=\varphi-\varphi_0\), their solution is
\begin{equation}
 \begin{pmatrix}a(\varphi)\\b(\varphi)\end{pmatrix}
 =\begin{pmatrix}
 \cos(c_0\Delta\varphi)&\sin(c_0\Delta\varphi)\\
 -\sin(c_0\Delta\varphi)&\cos(c_0\Delta\varphi)
 \end{pmatrix}
 \begin{pmatrix}a(\varphi_0)\\b(\varphi_0)\end{pmatrix}.
 \label{eq:latitude-transport}
\end{equation}
The quantity \(a^2+b^2=g(V,V)\) is preserved. The frame extends smoothly
around the latitude even though a single longitude chart does not. After
one circuit the frame returns to itself, but the vector has generally
rotated by \(-2\pi\cos\vartheta_0\), understood modulo \(2\pi\).
Thus transport around a closed path need not be the identity.

This calculation also separates two meanings of transport. The coordinate
fields commute, \([\partial_\varphi,\partial_\vartheta]=0\), but
\[
 \nabla_{\partial_\varphi}\partial_\vartheta
 =\cot\vartheta\,\partial_\varphi.
\]
Rotations preserve \(\partial_\vartheta\) on the angular patch, but that field
is not parallel along a general latitude. Lie invariance and parallelism
are distinct even for a metric-compatible, torsion-free connection.

\begin{studentwork}[2.5cm]{parallel transport on the sphere}
Recover \eqref{eq:latitude-transport} using coordinate components.
Determine the transport map after one circuit at
\(\vartheta_0=\pi/2\) and \(\vartheta_0=\pi/3\), and verify
preservation of the metric pairing.
\end{studentwork}

\subsection{Coordinates inertial at one event}
\label{sec:inertial-point}

The Levi--Civita connection permits a stronger coordinate choice than the
pointwise orthonormalization in Section~\ref{sec:local-lorentz}. First choose
coordinates \(x\) with \(x(p)=0\) and \(g_{\mu\nu}(p)=\eta_{\mu\nu}\).
A quadratic coordinate change can leave this basis unchanged at \(p\)
while adjusting the connection coefficients through its second derivatives.
Define new coordinates locally through
\[
 x^\rho=y^\rho-\frac12\Gamma^\rho{}_{\mu\nu}(p)y^\mu y^\nu.
\]
The Jacobian at the origin is the identity, so the inverse function theorem
makes this a valid local change of coordinates. Since the lower indices of
the Levi--Civita coefficients are symmetric, their transformation law gives
\[
 \widetilde\Gamma^\rho{}_{\mu\nu}(p)
 =\Gamma^\rho{}_{\mu\nu}(p)
   +\left.\frac{\partial^2x^\rho}{\partial y^\mu\partial y^\nu}\right|_p
 =0.
\]
The metric at the point remains \(\eta\), and metric compatibility then
implies
\[
 \widetilde g_{\mu\nu}(p)=\eta_{\mu\nu},\qquad
 \partial_\rho\widetilde g_{\mu\nu}(p)=0.
\]
These are local inertial conditions at a single event. They do not require
the second metric derivatives to vanish and do not establish flatness on a
neighborhood. The analogous Riemannian construction replaces \(\eta\) by
\(\delta\). The quadratic construction fixes only the coordinate data needed
at the point, not a unique chart away from it.

\section{Boost coordinates in Minkowski spacetime}
\label{sec:boost-app}

This application uses a noninertial coordinate chart on flat spacetime.
It combines the metric, a Killing field, covariant differentiation, and
proper time without introducing a gravitational field.

\subsection{The chart and its metric}

Set \(c=1\) and use inertial coordinates \((t,x,y,z)\). On the right-hand
wedge \(x>|t|\), introduce the rapidity coordinate \(\chi\) and radial
coordinate \(\rho>0\). Boosts preserve the hyperbolas
\(x^2-t^2=\rho^2\), and \(\chi\) parametrizes motion along each one.
Thus set
\begin{equation}
 t=\rho\sinh\chi,\qquad x=\rho\cosh\chi.
 \label{eq:boost-coordinates}
\end{equation}
The inverse is
\[
 \rho=\sqrt{x^2-t^2},\qquad
 \chi=\operatorname{artanh}(t/x).
\]
The Jacobian determinant of \((\chi,\rho)\mapsto(t,x)\) is \(\rho\), so it
is nonzero on this domain. Differentiation gives
\[
 \dd t=\rho\cosh\chi\,\dd\chi+\sinh\chi\,\dd\rho,\qquad
 \dd x=\rho\sinh\chi\,\dd\chi+\cosh\chi\,\dd\rho.
\]
Substituting in the Minkowski line element cancels the mixed terms,
\begin{equation}
 \dd s^2=\rho^2(\dd\chi)^2-(\dd\rho)^2-(\dd y)^2-(\dd z)^2.
 \label{eq:rindler-metric}
\end{equation}
Thus, in the order \((\chi,\rho,y,z)\),
\[
 (g_{\mu\nu})=\operatorname{diag}(\rho^2,-1,-1,-1),\qquad
 (g^{\mu\nu})=\operatorname{diag}(\rho^{-2},-1,-1,-1).
\]
The components vary, but the spacetime is still Minkowski spacetime because
this metric was obtained by an invertible change of chart. The chart does
not cover the wedge boundary. The inertial metric remains regular there.

\subsection{Boost symmetry and covariant differentiation}

The coordinate field
\[
 K=\partial_\chi=x\partial_t+t\partial_x
\]
generates active Lorentz boosts. Its flow is \(\chi\mapsto\chi+s\), with
\(\rho,y,z\) fixed. In the passive boost convention of SR1, the same
matrix is written \(L(-s)\). As a global field on Minkowski spacetime,
\[
 g(K,K)=x^2-t^2.
\]
It is timelike on the left and right wedges, null on their boundaries when
nonzero, and spacelike where \(|t|>|x|\). It vanishes when \(t=x=0\).
On the right wedge it is future directed for the standard time orientation.

The metric components in \eqref{eq:rindler-metric} do not depend on
\(\chi\), and the components of \(K=\partial_\chi\) are constant there.
The Lie-derivative formula therefore gives \(\Lie_Kg=0\).
The nonzero Levi--Civita coefficients are
\begin{equation}
 \Gamma^\rho{}_{\chi\chi}=\rho,\qquad
 \Gamma^\chi{}_{\chi\rho}=\Gamma^\chi{}_{\rho\chi}=\frac1\rho.
 \label{eq:rindler-gamma}
\end{equation}
For example,
\(\Gamma^\rho{}_{\chi\chi}
=\tfrac12(-1)(-\partial_\rho\rho^2)=\rho\).
Consequently,
\[
 [\partial_\chi,\partial_\rho]=0,
 \qquad \nabla_{\partial_\chi}\partial_\rho=\frac1\rho\partial_\chi.
\]
The coordinate frame commutes, but it is not a parallel frame.

\subsection{An accelerated clock}

Consider a future-directed worldline with \(\rho=\rho_0>0\) and constant
\(y,z\). Its proper time satisfies \(\dd\tau=\rho_0\,\dd\chi\), so
\[
 U=\frac{\dd\gamma}{\dd\tau}=\frac1{\rho_0}\partial_\chi,
 \qquad g(U,U)=1.
\]
The four-acceleration is the covariant derivative of the four-velocity,
\[
 a=\nabla_UU
   =\frac1{\rho_0^2}\nabla_{\partial_\chi}\partial_\chi
   =\frac1{\rho_0}\partial_\rho.
\]
It has invariant square
\[
 g(a,a)=-\frac1{\rho_0^2},\qquad g(a,U)=0.
\]
The orthogonality also follows by differentiating the constant
normalization,
\[
 0=U\bigl(g(U,U)\bigr)=2g(\nabla_UU,U)=2g(a,U).
\]
The observer has constant spatial coordinate labels, but nonzero
four-acceleration. Being at rest in this coordinate system is therefore
not the same as moving inertially.

In inertial coordinates, take \(\chi=\tau/\rho_0\) after choosing the
origins of both parameters. The component arrays in the inertial
coordinate basis are
\[
 (U^\mu)=(\cosh\chi,\sinh\chi,0,0),\qquad
 (a^\mu)=\frac1{\rho_0}(\sinh\chi,\cosh\chi,0,0).
\]
Ordinary proper-time differentiation is valid in this Cartesian inertial
frame because its connection coefficients vanish. It gives the same
acceleration and invariants. With length-valued \(\rho_0\) and \(c\) restored,
\[
 \dd\tau=\frac{\rho_0}{c}\,\dd\chi,\qquad
 \sqrt{-g(a,a)}=\frac{c^2}{\rho_0},
\]
where \(U=\dd\gamma/\dd\tau\) now has \(g(U,U)=c^2\).

\begin{studentwork}[4.3cm]{parallel transport along the accelerated worldline}
Along the fixed-\(\rho_0\) worldline, find the parallel field
\(V=A(\chi)\partial_\chi+B(\chi)\partial_\rho\) with
\(A(0)=A_0\), \(B(0)=B_0\), in units \(c=1\). Verify that
\(\rho_0^2A^2-B^2\) is constant and explain why the observer's
four-velocity is not parallel along the same curve.
\end{studentwork}

\end{document}